% ------------------------------------------------------------------------------
% main.tex — standalone compile wrapper for the codoc method section.
%
% For the actual ICLR submission, copy the preamble block marked [ICLR] into the
% ICLR template's main file and % ------------------------------------------------------------------------------
% method.tex — the codoc method/algorithm section.
%
% Self-contained: % ------------------------------------------------------------------------------
% method.tex — the codoc method/algorithm section.
%
% Self-contained: % ------------------------------------------------------------------------------
% method.tex — the codoc method/algorithm section.
%
% Self-contained: \input{method} from the ICLR template main file.
% Requires (all standard in the ICLR template ecosystem):
%   amsmath, amssymb, amsthm, booktabs, algorithm, algpseudocode
% plus the \newtheorem{definition}{Definition} declaration and the
% `problembox` environment + notation macros defined in main.tex.
% ------------------------------------------------------------------------------

\section{The \codoc{} System}
\label{sec:method}

\codoc{} maintains a human-authored \emph{feature tree} --- a hierarchy of named
units of intent --- bidirectionally synchronized with the code of a repository.
Unlike repository-level code graphs built \emph{from} code as retrieval context
for agents \citep{repograph2025,locagent2025,graphcoder2024,codexgraph2024}, the
tree is first-class authored intent: it is never regenerated from code, and the
attribution of code to features is a secondary, continuously repaired index.
Synchronization is asymmetric --- code changes are \emph{reflected} into the tree
(Loop~A, \S\ref{sec:loop-a}), while tree edits are \emph{realized} into code by a
coding agent (Loop~B, \S\ref{sec:loop-b}) --- and conflict resolution is
deterministic: pending document intent always wins over code-side observations
(\S\ref{sec:conflicts}).

\begin{problembox}{Problem statement}
Given a repository whose indexed chunk set is $\C$ and a human-authored feature
tree $\T = (\F, \prec)$, maintain (i)~a binding map
$\B \colon \C \to \F \cup \{\bot\}$ attributing every chunk to at most one
feature, and (ii)~the invariant that $\T$ expresses current human intent over
the code, under concurrent edits to either side by humans and autonomous
agents. Every mutation must be attributable (who, with what commitment, caused
by what), and divergence must resolve to a synchronization state
$\sigma = \synced$ without losing authored intent.
\end{problembox}

\subsection{Preliminaries}
\label{sec:prelim}

\begin{definition}[Chunk set]
\label{def:chunks}
The repository is indexed into a set of AST-level chunks
$\C = \{c_1, \dots, c_n\}$, where each chunk
$c = (\mathit{file},\, \mathit{sym},\, h^{\mathrm{tok}},\, h^{\mathrm{ast}})$
carries its file path, its qualified symbol path within the file, a
\emph{content hash} $h^{\mathrm{tok}}$ over the comment- and
whitespace-normalized token stream, and an \emph{AST-shape hash}
$h^{\mathrm{ast}}$ over the token-erased tree structure. The pair
$(\mathit{file}, \mathit{sym})$ is the chunk's identity key; the two hashes are
its relocation signals (Definition~\ref{def:reloc}).
\end{definition}

\begin{definition}[Feature tree]
\label{def:tree}
The feature tree is $\T = (\F, \prec)$, where $\F$ is a finite set of features
$f = (\mathit{id},\, \mathit{title},\, \mathit{desc},\, \mathit{realized})$ and
$\prec\, \subseteq \F \times \F$ is a forest-shaped parent relation. Titles and
descriptions are free prose authored by humans (or proposed by agents and
accepted by humans); $\mathit{realized} \in \{0,1\}$ marks whether the feature
is backed by code yet ($\mathit{realized}=0$ denotes an accepted plan
placeholder, flipped to $1$ by the first binding).
\end{definition}

\begin{definition}[Binding map]
\label{def:binding}
Attribution is a partial function $\B \colon \C \to \F \cup \{\bot\}$: each
chunk binds to at most one feature (enforced as a uniqueness constraint on the
chunk identity key), while one feature typically owns many chunks across many
files. Each binding stores the bound chunk's $h^{\mathrm{tok}}$ as its
\emph{fingerprint}, making staleness locally detectable. We write
$\B^{-1}(f) \subseteq \C$ for the chunks a feature owns and define coverage
$\mathrm{cov}(\B) = \lvert\{c : \B(c) \neq \bot\}\rvert / \lvert\C\rvert$.
\end{definition}

\begin{definition}[Code graph]
\label{def:graph}
Following the heterogeneous-graph convention of repository-level code
representations \citep{locagent2025,graphcoder2024}, the dependency structure
of the code is $\G = (\V, \E, \A, \R)$ with node-type map
$\tau \colon \V \to \A$, $\A = \{\textit{file}, \textit{class},
\textit{function}\}$, and edge-type map $\varphi \colon \E \to \R$,
$\R = \{\textit{contain}, \textit{call}, \textit{import}, \textit{inherit}\}$.
$\G$ is a derived cache, rebuildable from $\C$ at any time. The binding map
lifts $\G$ to a feature-coupling graph and defines the placement heuristic
\begin{equation}
\label{eq:neighbor}
\mathrm{nbf}(c) \;=\; \argmax_{f \in \F}\;
\bigl\lvert \{\, (v, v') \in \E \;:\; v' = \mathrm{sym}(c),\;
\B(\mathrm{chunk}(v)) = f \,\} \bigr\rvert ,
\end{equation}
the feature owning the most call/import edges incident to a new chunk $c$ ---
the deterministic fallback of Loop~A's coverage net
(Algorithm~\ref{alg:loop-a}, lines~\ref{line:covnet-start}--\ref{line:covnet-end}).
\end{definition}

\begin{definition}[Changeset]
\label{def:changeset}
A pass over the index yields $\Delta = (\Delta^{+}, \Delta^{-}, \Delta^{\sim})$:
chunks \emph{added} (present, unbound), \emph{removed} (a binding whose key has
left the index), and \emph{modified} (bound, but
$h^{\mathrm{tok}} \neq$ fingerprint). \codoc{} computes $\Delta$ in two modes:
a \emph{temporal} diff of consecutive index snapshots (cheap, runs on every
save) and a \emph{state-based} reconciliation that compares the index directly
against $\B$ (idempotent and self-healing; the authority pass). Only the latter
may retire features or amend prose, so destructive judgments are never made
from a transient mid-edit view.
\end{definition}

\begin{definition}[Relocation]
\label{def:reloc}
For $(r, a) \in \Delta^{-} \times \Delta^{+}$: $a$ is a \emph{move} of $r$ iff
$h^{\mathrm{tok}}_a = h^{\mathrm{tok}}_r$ (identical content at a new key), and
a \emph{rename} of $r$ iff $h^{\mathrm{ast}}_a = h^{\mathrm{ast}}_r$ uniquely
within the same file (same shape, new name; $1{:}1$ matches only). Relocations
re-attach attribution deterministically --- no model call, no risk of dropped
or duplicated attribution.
\end{definition}

\begin{definition}[Change ledger]
\label{def:ledger}
Every mutation of $(\T, \B)$ is an event in an append-only ledger $L$. A ledger
entry $\ell = (\mathit{op},\, f,\, \alpha,\, \mu,\, \delta)$ records the
operation, its target, the \emph{actor}
$\alpha \in \{\textsf{human}\} \cup \mathit{Agents} \cup \{\textsf{loop}\}$,
the \emph{mode} $\mu \in \{\pen, \suggest, \auto\}$ (committed edit, reviewable
proposal, or mechanical maintenance), and a \emph{causality link} $\delta$ ---
the directive, suggestion, or event id this change implements ($\bot$ if
spontaneous). A \emph{proposal} is an unapplied event; accepting it applies the
operation and consumes the event.
\end{definition}

\begin{definition}[Hold set and synchronization state]
\label{def:hold}
$H \subseteq \F$ is the set of features with pending document-ahead intent: a
live suggestion not yet drained, or a queued directive not yet implemented.
The synchronization state is
$\sigma \in \{\synced, \codedrift, \treedirty, \awaiting, \realizing\}$,
derived (not stored): pending code-side proposals put $\sigma = \codedrift$, a
non-empty directive queue forces $\sigma = \awaiting$ (a later code-side pass
cannot silently clobber queued work), and an active implementation pass reports
$\sigma = \realizing$.
\end{definition}

Table~\ref{tab:notation} summarizes the notation.

\begin{table}[t]
\centering
\caption{Notation.}
\label{tab:notation}
\small
\begin{tabular}{@{}ll@{}}
\toprule
Symbol & Meaning \\
\midrule
$\C$ & chunk set; $c = (\mathit{file}, \mathit{sym}, h^{\mathrm{tok}}, h^{\mathrm{ast}})$ \\
$\T = (\F, \prec)$ & feature tree (features + forest parent relation) \\
$\B \colon \C \to \F \cup \{\bot\}$ & binding map (each chunk owned by $\leq 1$ feature) \\
$\G = (\V, \E, \A, \R)$ & heterogeneous code dependency graph (derived cache) \\
$\Delta = (\Delta^{+}, \Delta^{-}, \Delta^{\sim})$ & changeset: added / removed / modified chunks \\
$h^{\mathrm{tok}}, h^{\mathrm{ast}}$ & content hash (move-invariant), AST-shape hash (rename-invariant) \\
$L$; $\ell = (\mathit{op}, f, \alpha, \mu, \delta)$ & append-only ledger; entry = op, target, actor, mode, cause \\
$H \subseteq \F$ & hold set (features with pending doc-ahead intent) \\
$\sigma$ & synchronization state \\
$\mathrm{nbf}(c)$ & graph-neighbor feature of a chunk (Eq.~\ref{eq:neighbor}) \\
$D$; $d$ & directive queue; a single realization directive \\
\bottomrule
\end{tabular}
\end{table}

\subsection{Loop A: Reflecting Code into the Tree}
\label{sec:loop-a}

Loop~A (Algorithm~\ref{alg:loop-a}) closes the gap from code to intent. Its
design principle is \emph{determinism first}: every change that can be resolved
mechanically --- refreshing fingerprints, detaching deleted chunks, carrying
attribution across moves and renames --- is resolved without a model
(lines~\ref{line:auto-start}--\ref{line:reloc-end}). A single LLM pass
(line~\ref{line:llm}) handles only the judgment calls that remain: placing
genuinely new code, repairing features that lost their last binding, and (in
the authority pass) refreshing prose whose code changed underneath it. The
pass sees the residual changeset, a file-locality subtree of $\T$, and
\emph{every} feature title; supplying the full title set to one call is what
prevents duplicate features without any post-hoc deduplication machinery.

Model output is not trusted with destructive authority. Proposed operations
are split by safety: binding maintenance and small description amendments
(bounded by an edit-distance ratio) auto-apply with $\mu = \auto$; structural
operations (add / move / retire) become unapplied proposal events with
$\mu = \suggest$ that a human accepts or rejects. Operations targeting held
features are dropped (line~\ref{line:hold}; \S\ref{sec:conflicts}), and a
\emph{coverage net}
(lines~\ref{line:covnet-start}--\ref{line:covnet-end}) guarantees no added
chunk is silently dropped: anything the model failed to place attaches to its
graph-neighbor feature $\mathrm{nbf}(c)$ or surfaces as an explicit
\textsf{add} proposal. Finally, stale-proposal garbage collection keeps the
pass idempotent: a pending \textsf{add} whose chunks have since bound
elsewhere, or a pending \textsf{retire} whose feature regained code, is
dropped, so repeated passes converge to $\sigma = \synced$ instead of
accumulating false positives.

\begin{algorithm}[t]
\caption{Loop A --- code $\to$ tree reflection (one pass)}
\label{alg:loop-a}
\begin{algorithmic}[1]
\Require index snapshot of $\C$; store $(\T, \B, L)$; hold set $H$; authority flag $\mathit{auth}$
\Ensure updated $(\T, \B, L)$; refreshed $\sigma$
\State $\Delta \gets \textsc{Changeset}(\C, \B)$ \Comment{temporal diff, or state-based iff $\mathit{auth}$}
\State \textsc{GcStaleProposals}$(L, \Delta)$ \Comment{drop superseded \textsf{add}/\textsf{retire} proposals}
\For{$c \in \Delta^{\sim}$} \textsc{Apply}(\textsf{refresh}$(c)$, $\mu{=}\auto$) \EndFor \label{line:auto-start}
\For{$c \in \Delta^{-}$} \textsc{Apply}(\textsf{detach}$(c)$, $\mu{=}\auto$) \EndFor
\For{$(r, a) \in \textsc{Relocations}(\Delta^{-}, \Delta^{+})$} \Comment{Def.~\ref{def:reloc}: $h^{\mathrm{tok}}$ move / $h^{\mathrm{ast}}$ rename}
    \State \textsc{Apply}(\textsf{attach}$(a, \B(r))$, $\mu{=}\auto$) \label{line:reloc-end}
\EndFor
\State $U \gets \{c \in \Delta^{+} : \B(c) = \bot\}$;\quad
       $E \gets \{f : \B^{-1}(f) = \emptyset \wedge f \notin H\}$ \Comment{unbound adds; emptied features}
\If{$U \neq \emptyset \,\vee\, E \neq \emptyset \,\vee\, (\mathit{auth} \wedge \textsc{StaleProse}(\Delta^{\sim}, H))$}
    \State $\mathit{ops} \gets \textsc{LLM}\bigl(\Delta,\ \textsc{Subtree}(\T, \Delta),\ \textsc{Titles}(\T),\ \G\bigr)$ \Comment{one call per pass} \label{line:llm}
    \State $\mathit{ops} \gets \{o \in \mathit{ops} : \neg\,\textsc{SuppressedByHold}(o, H)\}$ \label{line:hold}
    \For{$o \in \mathit{ops}$}
        \If{\textsc{Safe}$(o)$} \textsc{Apply}$(o, \mu{=}\auto)$
        \Else\ \textsc{Propose}$(o, \mu{=}\suggest)$ \Comment{unapplied event; human verdict pending}
        \EndIf
    \EndFor
\EndIf
\For{$c \in U$ with $\B(c) = \bot$ still} \label{line:covnet-start} \Comment{coverage net: nothing is dropped}
    \If{$\mathrm{nbf}(c) \neq \bot$} \textsc{Apply}(\textsf{attach}$(c, \mathrm{nbf}(c))$, $\mu{=}\auto$)
    \Else\ \textsc{Propose}(\textsf{add}$(c)$, $\mu{=}\suggest$)
    \EndIf
\EndFor \label{line:covnet-end}
\State \Return $(\T, \B, L)$, \textsc{RefreshState}$()$
\end{algorithmic}
\end{algorithm}

Two authority levels share this pipeline. The \emph{temporal} pass (run on
every save) computes $\Delta$ from consecutive snapshots and may never retire a
feature or amend prose --- a feature that looks emptied mid-edit usually
rebinds on the next save. The \emph{authority} pass (run at agent-session
close, on daemon startup, and from explicit sync) computes $\Delta$ state-based
against $\B$, is idempotent, recovers missed cycles, and is the only pass
allowed to raise \textsf{retire} and stale-prose \textsf{amend} proposals.

\subsection{Loop B: Realizing Tree Edits into Code}
\label{sec:loop-b}

Loop~B (Algorithm~\ref{alg:loop-b}) closes the opposite gap. Edits to the
rendered tree document are parsed and diffed against the store; verdicts on
pending proposals are drained from the review inbox; and live document-ahead
suggestions are applied by the loop itself (the human's only verb on their own
suggestion is \emph{withdraw} --- application is the agent side's
acknowledgment, stamped $\mu = \suggest$ with $\delta$ = the suggestion id).

Two ordering invariants keep text and store coherent
(lines~\ref{line:snapshot} and~\ref{line:rerender}). The text diff is
snapshotted \emph{before} any store mutation: verdict accepts move the store
ahead of the on-disk text, and diffing afterwards would misread the stale text
as a human edit that reverts the accepted change. Symmetrically, any pass that
mutated the store re-renders the text before returning, so the next pass diffs
a caught-up document.

Not every tree edit means ``write code.'' The \emph{imperative gate}
(line~\ref{line:imperative}) classifies each applied edit: descriptive prose
(``validates the token'') merely persists, while imperative prose (``should
validate the token'', ``Add retry logic''), an accepted unrealized plan node,
or retiring a feature that owns code each mint a \emph{directive} $d$ --- a
natural-language work order assembled from the feature's description and its
bound symbols, scoped to the files the feature owns. Directives are queued for
the user's \emph{live} coding-agent session rather than a headless model: the
session implements each directive, reflects the new code back through the
agent-side reflection API with $\delta = d$, and deletes the queue. Loop~A
then closes the cycle, with every resulting ledger entry carrying the causal
chain from the original document edit.

Verdict semantics are deliberately asymmetric for the destructive case
(line~\ref{line:retire}): accepting a machine-raised \textsf{retire} from the
inbox is \emph{detach-only} --- the feature is hidden and its chunks freed for
re-attribution, but no code-deletion directive is ever queued, because a retire
raised from transient drift may be a false positive and deleting code on a
single click is the most destructive failure mode. Only an explicit human
retire edit in the document text (or an agent proposal explicitly flagged to
delete code) queues removal work.

\begin{algorithm}[t]
\caption{Loop B --- tree $\to$ code realization (one pass)}
\label{alg:loop-b}
\begin{algorithmic}[1]
\Require tree text $T_{\mathrm{txt}}$; inbox verdicts $V$; intents $I$; store $(\T, \B, L)$
\Ensure updated $(\T, \B, L)$; directive queue $D$; refreshed $\sigma$
\State $\mathit{ops}_{\mathrm{user}} \gets \textsc{Diff}(T_{\mathrm{txt}}, \T)$ \Comment{snapshot BEFORE mutations (invariant 1)} \label{line:snapshot}
\State $D \gets \emptyset$
\For{$(e, \mathit{accept}) \in V$} \Comment{drain human verdicts on pending proposals}
    \If{$\mathit{accept}$}
        \State \textsc{Apply}$(e.\mathit{op}, \alpha{=}\textsf{human}, \mu{=}\pen)$; \textsc{Consume}$(e)$
        \If{$e.\mathit{op}$ is \textsf{retire} $\wedge\, \neg e.\mathit{op}.\mathit{deleteCode}$} \label{line:retire}
            \State \textsc{DetachAll}$(\B^{-1}(e.\mathit{op}.f))$ \Comment{detach-only: never queue code deletion}
        \ElsIf{\textsc{ImpliesCode}$(e.\mathit{op})$} $D \gets D \cup \{\textsc{Directive}(e.\mathit{op},\, \delta{=}e)\}$
        \EndIf
    \Else\ \textsc{Consume}$(e)$ \Comment{reject: delete the proposal}
    \EndIf
\EndFor
\For{$o \in \mathit{ops}_{\mathrm{user}}$} \Comment{direct edits are intentional: apply immediately}
    \State \textsc{Apply}$(o, \alpha, \mu \text{ from annotation channel (default \textsf{human}/\pen)})$
    \If{\textsc{ImpliesCode}$(o)$} $D \gets D \cup \{\textsc{Directive}(o, \delta{=}\ell_o)\}$ \label{line:imperative}
    \EndIf
\EndFor
\For{$i \in I$ with payload, fresh, unsatisfied} \Comment{doc-ahead suggestions: the loop applies them}
    \State \textsc{Apply}(\textsf{amend}$(i)$, $\alpha{=}i.\alpha$, $\mu{=}\suggest$, $\delta{=}i.\mathit{id}$)
    \If{\textsc{ImpliesCode}$(i)$} $D \gets D \cup \{\textsc{Directive}(i, \delta{=}i.\mathit{id})\}$
    \EndIf
\EndFor
\If{store mutated} \textsc{Render}$(\T) \to T_{\mathrm{txt}}$ \Comment{invariant 2: text catches up} \label{line:rerender}
\EndIf
\If{$D \neq \emptyset$}
    \State mint ids for $D$; queue $D$ for the live agent session; $H \gets H \cup \{d.f : d \in D\}$
    \State $\sigma \gets \awaiting$ \Comment{the agent implements, reflects with $\delta = d$, clears the queue}
\Else
    \State $\sigma \gets \textsc{RefreshState}()$
\EndIf
\State \Return $(\T, \B, L)$, $D$, $\sigma$
\end{algorithmic}
\end{algorithm}

\subsection{Conflict Resolution: the Ledger and Doc-Wins Holds}
\label{sec:conflicts}

Both loops route every decision through one classification table
(Table~\ref{tab:classify}): each detected change, from either side and by any
actor, maps to exactly one reaction. Three properties follow.

\paragraph{Provenance.} Every mutation is a ledger entry
(Definition~\ref{def:ledger}). The actor/mode pair separates \emph{who} changed
something from \emph{how committed} the change is, which drives both review
routing (only $\mu = \suggest$ items need verdicts) and the interface's
authorship rendering; $\delta$ threads causality across the realize cycle, so
a code change surfaced back by Loop~A is grouped under the document edit that
caused it.

\paragraph{Doc-wins holds.} Conflict resolution is a deterministic priority
rule rather than a merge heuristic. While $f \in H$ (a suggestion or queued
directive is pending on $f$), code-side \emph{intent} operations on $f$ ---
\textsf{amend}, \textsf{retire}, \textsf{move} --- are suppressed:
\begin{equation}
\label{eq:docwins}
\mathrm{suppress}(o, H) \;\iff\;
o.\mathit{kind} \in \{\textsf{amend}, \textsf{retire}, \textsf{move}\}
\;\wedge\; o.f \in H .
\end{equation}
Binding maintenance (\textsf{attach}/\textsf{detach}/\textsf{refresh}) is never
suppressed --- bindings are attribution, not intent, and must stay correct
regardless of who is mid-edit. Holds release when the suggestion is drained or
the directive queue completes, with a staleness backstop so an abandoned
suggestion cannot hold a feature indefinitely. The rule is the
tree-is-authoritative analogue of last-writer-wins: the side that expresses
\emph{intent} (the document) always outranks the side that merely
\emph{observes} (the code reflector).

\paragraph{Verdict grammar.} Who resolves what is structural:
accept/reject is exclusively a \emph{human} verb on \emph{code-ahead} items
(machine proposals), while a \emph{doc-ahead} suggestion is applied by the
machine side and offers its human author only \emph{withdraw}. No state is
double-resolvable, which eliminates accept/revert races by construction.

\begin{table}[t]
\centering
\caption{The change-classification table. Every change, from either side and by
any actor, maps to one reaction. $\auto$ reactions apply immediately;
$\suggest$ reactions await a human verdict; \textbf{D} marks reactions that
mint a realization directive.}
\label{tab:classify}
\small
\setlength{\tabcolsep}{4.5pt}
\begin{tabular}{@{}cllll@{}}
\toprule
\# & Side & Detected change & Origin & Reaction \\
\midrule
1 & code & modified bound chunk & any & \textsf{refresh} binding ($\auto$) \\
2 & code & removed bound chunk & any & \textsf{detach} ($\auto$); emptied $\wedge\, f \notin H \Rightarrow$ \textsf{retire} ($\suggest$, auth.\ pass) \\
3 & code & relocated chunk (Def.~\ref{def:reloc}) & any & deterministic \textsf{attach} to prior feature ($\auto$) \\
4 & code & added unbound chunk & any & LLM placement ($\auto$/$\suggest$); fallback $\mathrm{nbf}$ / \textsf{add} ($\suggest$) \\
5 & code & in-place modify, realized + prose & any & stale-prose \textsf{amend} (small $\to \auto$, large $\to \suggest$) \\
6 & code & any change during realize epoch & agent & rows 1--5 with $\delta = $ directive id \\
\midrule
7 & doc & descriptive amend / title edit & human, $\pen$ & apply immediately; no directive \\
8 & doc & imperative amend / plan add / retire-with-code & human, $\pen$ & apply + mint directive (\textbf{D}) \\
9 & doc & any edit & human, $\suggest$ & pending intent (hold); applied by Loop B drain $\to$ rows 7/8 \\
10 & doc & structural op via reflection API & agent & pending proposal ($\suggest$) \\
11 & doc & safe op via reflection API & agent & auto-apply; ledger $\alpha = \textsf{agent}$, $\mu = \auto$ \\
\midrule
12 & verdict & accept / reject & human & apply+consume / delete (\textsf{retire} accept is detach-only) \\
13 & code & drift on held feature ($f \in H$) & any & bindings maintained; \textsf{amend}/\textsf{retire}/\textsf{move} suppressed (Eq.~\ref{eq:docwins}) \\
\bottomrule
\end{tabular}
\end{table}

\subsection{Bootstrap}
\label{sec:bootstrap}

The initial tree is built by two scoped passes rather than one global
clustering. A \emph{per-file} pass issues one model call per source file,
seeing only that file's chunks and their incident graph edges --- making
cross-file ``junk drawer'' features structurally impossible --- and a
deterministic coverage step folds any chunk the model missed into the file's
largest feature. An \emph{organization} pass then groups the resulting
file-level features under a small number of thematic parents, using
feature-to-feature coupling aggregated from $\G$ as evidence. Both passes emit
ordinary ledger events through the same application seam as the loops, so
bootstrap output is curated (and corrected) with exactly the same verdict
machinery as any later change.
 from the ICLR template main file.
% Requires (all standard in the ICLR template ecosystem):
%   amsmath, amssymb, amsthm, booktabs, algorithm, algpseudocode
% plus the \newtheorem{definition}{Definition} declaration and the
% `problembox` environment + notation macros defined in main.tex.
% ------------------------------------------------------------------------------

\section{The \codoc{} System}
\label{sec:method}

\codoc{} maintains a human-authored \emph{feature tree} --- a hierarchy of named
units of intent --- bidirectionally synchronized with the code of a repository.
Unlike repository-level code graphs built \emph{from} code as retrieval context
for agents \citep{repograph2025,locagent2025,graphcoder2024,codexgraph2024}, the
tree is first-class authored intent: it is never regenerated from code, and the
attribution of code to features is a secondary, continuously repaired index.
Synchronization is asymmetric --- code changes are \emph{reflected} into the tree
(Loop~A, \S\ref{sec:loop-a}), while tree edits are \emph{realized} into code by a
coding agent (Loop~B, \S\ref{sec:loop-b}) --- and conflict resolution is
deterministic: pending document intent always wins over code-side observations
(\S\ref{sec:conflicts}).

\begin{problembox}{Problem statement}
Given a repository whose indexed chunk set is $\C$ and a human-authored feature
tree $\T = (\F, \prec)$, maintain (i)~a binding map
$\B \colon \C \to \F \cup \{\bot\}$ attributing every chunk to at most one
feature, and (ii)~the invariant that $\T$ expresses current human intent over
the code, under concurrent edits to either side by humans and autonomous
agents. Every mutation must be attributable (who, with what commitment, caused
by what), and divergence must resolve to a synchronization state
$\sigma = \synced$ without losing authored intent.
\end{problembox}

\subsection{Preliminaries}
\label{sec:prelim}

\begin{definition}[Chunk set]
\label{def:chunks}
The repository is indexed into a set of AST-level chunks
$\C = \{c_1, \dots, c_n\}$, where each chunk
$c = (\mathit{file},\, \mathit{sym},\, h^{\mathrm{tok}},\, h^{\mathrm{ast}})$
carries its file path, its qualified symbol path within the file, a
\emph{content hash} $h^{\mathrm{tok}}$ over the comment- and
whitespace-normalized token stream, and an \emph{AST-shape hash}
$h^{\mathrm{ast}}$ over the token-erased tree structure. The pair
$(\mathit{file}, \mathit{sym})$ is the chunk's identity key; the two hashes are
its relocation signals (Definition~\ref{def:reloc}).
\end{definition}

\begin{definition}[Feature tree]
\label{def:tree}
The feature tree is $\T = (\F, \prec)$, where $\F$ is a finite set of features
$f = (\mathit{id},\, \mathit{title},\, \mathit{desc},\, \mathit{realized})$ and
$\prec\, \subseteq \F \times \F$ is a forest-shaped parent relation. Titles and
descriptions are free prose authored by humans (or proposed by agents and
accepted by humans); $\mathit{realized} \in \{0,1\}$ marks whether the feature
is backed by code yet ($\mathit{realized}=0$ denotes an accepted plan
placeholder, flipped to $1$ by the first binding).
\end{definition}

\begin{definition}[Binding map]
\label{def:binding}
Attribution is a partial function $\B \colon \C \to \F \cup \{\bot\}$: each
chunk binds to at most one feature (enforced as a uniqueness constraint on the
chunk identity key), while one feature typically owns many chunks across many
files. Each binding stores the bound chunk's $h^{\mathrm{tok}}$ as its
\emph{fingerprint}, making staleness locally detectable. We write
$\B^{-1}(f) \subseteq \C$ for the chunks a feature owns and define coverage
$\mathrm{cov}(\B) = \lvert\{c : \B(c) \neq \bot\}\rvert / \lvert\C\rvert$.
\end{definition}

\begin{definition}[Code graph]
\label{def:graph}
Following the heterogeneous-graph convention of repository-level code
representations \citep{locagent2025,graphcoder2024}, the dependency structure
of the code is $\G = (\V, \E, \A, \R)$ with node-type map
$\tau \colon \V \to \A$, $\A = \{\textit{file}, \textit{class},
\textit{function}\}$, and edge-type map $\varphi \colon \E \to \R$,
$\R = \{\textit{contain}, \textit{call}, \textit{import}, \textit{inherit}\}$.
$\G$ is a derived cache, rebuildable from $\C$ at any time. The binding map
lifts $\G$ to a feature-coupling graph and defines the placement heuristic
\begin{equation}
\label{eq:neighbor}
\mathrm{nbf}(c) \;=\; \argmax_{f \in \F}\;
\bigl\lvert \{\, (v, v') \in \E \;:\; v' = \mathrm{sym}(c),\;
\B(\mathrm{chunk}(v)) = f \,\} \bigr\rvert ,
\end{equation}
the feature owning the most call/import edges incident to a new chunk $c$ ---
the deterministic fallback of Loop~A's coverage net
(Algorithm~\ref{alg:loop-a}, lines~\ref{line:covnet-start}--\ref{line:covnet-end}).
\end{definition}

\begin{definition}[Changeset]
\label{def:changeset}
A pass over the index yields $\Delta = (\Delta^{+}, \Delta^{-}, \Delta^{\sim})$:
chunks \emph{added} (present, unbound), \emph{removed} (a binding whose key has
left the index), and \emph{modified} (bound, but
$h^{\mathrm{tok}} \neq$ fingerprint). \codoc{} computes $\Delta$ in two modes:
a \emph{temporal} diff of consecutive index snapshots (cheap, runs on every
save) and a \emph{state-based} reconciliation that compares the index directly
against $\B$ (idempotent and self-healing; the authority pass). Only the latter
may retire features or amend prose, so destructive judgments are never made
from a transient mid-edit view.
\end{definition}

\begin{definition}[Relocation]
\label{def:reloc}
For $(r, a) \in \Delta^{-} \times \Delta^{+}$: $a$ is a \emph{move} of $r$ iff
$h^{\mathrm{tok}}_a = h^{\mathrm{tok}}_r$ (identical content at a new key), and
a \emph{rename} of $r$ iff $h^{\mathrm{ast}}_a = h^{\mathrm{ast}}_r$ uniquely
within the same file (same shape, new name; $1{:}1$ matches only). Relocations
re-attach attribution deterministically --- no model call, no risk of dropped
or duplicated attribution.
\end{definition}

\begin{definition}[Change ledger]
\label{def:ledger}
Every mutation of $(\T, \B)$ is an event in an append-only ledger $L$. A ledger
entry $\ell = (\mathit{op},\, f,\, \alpha,\, \mu,\, \delta)$ records the
operation, its target, the \emph{actor}
$\alpha \in \{\textsf{human}\} \cup \mathit{Agents} \cup \{\textsf{loop}\}$,
the \emph{mode} $\mu \in \{\pen, \suggest, \auto\}$ (committed edit, reviewable
proposal, or mechanical maintenance), and a \emph{causality link} $\delta$ ---
the directive, suggestion, or event id this change implements ($\bot$ if
spontaneous). A \emph{proposal} is an unapplied event; accepting it applies the
operation and consumes the event.
\end{definition}

\begin{definition}[Hold set and synchronization state]
\label{def:hold}
$H \subseteq \F$ is the set of features with pending document-ahead intent: a
live suggestion not yet drained, or a queued directive not yet implemented.
The synchronization state is
$\sigma \in \{\synced, \codedrift, \treedirty, \awaiting, \realizing\}$,
derived (not stored): pending code-side proposals put $\sigma = \codedrift$, a
non-empty directive queue forces $\sigma = \awaiting$ (a later code-side pass
cannot silently clobber queued work), and an active implementation pass reports
$\sigma = \realizing$.
\end{definition}

Table~\ref{tab:notation} summarizes the notation.

\begin{table}[t]
\centering
\caption{Notation.}
\label{tab:notation}
\small
\begin{tabular}{@{}ll@{}}
\toprule
Symbol & Meaning \\
\midrule
$\C$ & chunk set; $c = (\mathit{file}, \mathit{sym}, h^{\mathrm{tok}}, h^{\mathrm{ast}})$ \\
$\T = (\F, \prec)$ & feature tree (features + forest parent relation) \\
$\B \colon \C \to \F \cup \{\bot\}$ & binding map (each chunk owned by $\leq 1$ feature) \\
$\G = (\V, \E, \A, \R)$ & heterogeneous code dependency graph (derived cache) \\
$\Delta = (\Delta^{+}, \Delta^{-}, \Delta^{\sim})$ & changeset: added / removed / modified chunks \\
$h^{\mathrm{tok}}, h^{\mathrm{ast}}$ & content hash (move-invariant), AST-shape hash (rename-invariant) \\
$L$; $\ell = (\mathit{op}, f, \alpha, \mu, \delta)$ & append-only ledger; entry = op, target, actor, mode, cause \\
$H \subseteq \F$ & hold set (features with pending doc-ahead intent) \\
$\sigma$ & synchronization state \\
$\mathrm{nbf}(c)$ & graph-neighbor feature of a chunk (Eq.~\ref{eq:neighbor}) \\
$D$; $d$ & directive queue; a single realization directive \\
\bottomrule
\end{tabular}
\end{table}

\subsection{Loop A: Reflecting Code into the Tree}
\label{sec:loop-a}

Loop~A (Algorithm~\ref{alg:loop-a}) closes the gap from code to intent. Its
design principle is \emph{determinism first}: every change that can be resolved
mechanically --- refreshing fingerprints, detaching deleted chunks, carrying
attribution across moves and renames --- is resolved without a model
(lines~\ref{line:auto-start}--\ref{line:reloc-end}). A single LLM pass
(line~\ref{line:llm}) handles only the judgment calls that remain: placing
genuinely new code, repairing features that lost their last binding, and (in
the authority pass) refreshing prose whose code changed underneath it. The
pass sees the residual changeset, a file-locality subtree of $\T$, and
\emph{every} feature title; supplying the full title set to one call is what
prevents duplicate features without any post-hoc deduplication machinery.

Model output is not trusted with destructive authority. Proposed operations
are split by safety: binding maintenance and small description amendments
(bounded by an edit-distance ratio) auto-apply with $\mu = \auto$; structural
operations (add / move / retire) become unapplied proposal events with
$\mu = \suggest$ that a human accepts or rejects. Operations targeting held
features are dropped (line~\ref{line:hold}; \S\ref{sec:conflicts}), and a
\emph{coverage net}
(lines~\ref{line:covnet-start}--\ref{line:covnet-end}) guarantees no added
chunk is silently dropped: anything the model failed to place attaches to its
graph-neighbor feature $\mathrm{nbf}(c)$ or surfaces as an explicit
\textsf{add} proposal. Finally, stale-proposal garbage collection keeps the
pass idempotent: a pending \textsf{add} whose chunks have since bound
elsewhere, or a pending \textsf{retire} whose feature regained code, is
dropped, so repeated passes converge to $\sigma = \synced$ instead of
accumulating false positives.

\begin{algorithm}[t]
\caption{Loop A --- code $\to$ tree reflection (one pass)}
\label{alg:loop-a}
\begin{algorithmic}[1]
\Require index snapshot of $\C$; store $(\T, \B, L)$; hold set $H$; authority flag $\mathit{auth}$
\Ensure updated $(\T, \B, L)$; refreshed $\sigma$
\State $\Delta \gets \textsc{Changeset}(\C, \B)$ \Comment{temporal diff, or state-based iff $\mathit{auth}$}
\State \textsc{GcStaleProposals}$(L, \Delta)$ \Comment{drop superseded \textsf{add}/\textsf{retire} proposals}
\For{$c \in \Delta^{\sim}$} \textsc{Apply}(\textsf{refresh}$(c)$, $\mu{=}\auto$) \EndFor \label{line:auto-start}
\For{$c \in \Delta^{-}$} \textsc{Apply}(\textsf{detach}$(c)$, $\mu{=}\auto$) \EndFor
\For{$(r, a) \in \textsc{Relocations}(\Delta^{-}, \Delta^{+})$} \Comment{Def.~\ref{def:reloc}: $h^{\mathrm{tok}}$ move / $h^{\mathrm{ast}}$ rename}
    \State \textsc{Apply}(\textsf{attach}$(a, \B(r))$, $\mu{=}\auto$) \label{line:reloc-end}
\EndFor
\State $U \gets \{c \in \Delta^{+} : \B(c) = \bot\}$;\quad
       $E \gets \{f : \B^{-1}(f) = \emptyset \wedge f \notin H\}$ \Comment{unbound adds; emptied features}
\If{$U \neq \emptyset \,\vee\, E \neq \emptyset \,\vee\, (\mathit{auth} \wedge \textsc{StaleProse}(\Delta^{\sim}, H))$}
    \State $\mathit{ops} \gets \textsc{LLM}\bigl(\Delta,\ \textsc{Subtree}(\T, \Delta),\ \textsc{Titles}(\T),\ \G\bigr)$ \Comment{one call per pass} \label{line:llm}
    \State $\mathit{ops} \gets \{o \in \mathit{ops} : \neg\,\textsc{SuppressedByHold}(o, H)\}$ \label{line:hold}
    \For{$o \in \mathit{ops}$}
        \If{\textsc{Safe}$(o)$} \textsc{Apply}$(o, \mu{=}\auto)$
        \Else\ \textsc{Propose}$(o, \mu{=}\suggest)$ \Comment{unapplied event; human verdict pending}
        \EndIf
    \EndFor
\EndIf
\For{$c \in U$ with $\B(c) = \bot$ still} \label{line:covnet-start} \Comment{coverage net: nothing is dropped}
    \If{$\mathrm{nbf}(c) \neq \bot$} \textsc{Apply}(\textsf{attach}$(c, \mathrm{nbf}(c))$, $\mu{=}\auto$)
    \Else\ \textsc{Propose}(\textsf{add}$(c)$, $\mu{=}\suggest$)
    \EndIf
\EndFor \label{line:covnet-end}
\State \Return $(\T, \B, L)$, \textsc{RefreshState}$()$
\end{algorithmic}
\end{algorithm}

Two authority levels share this pipeline. The \emph{temporal} pass (run on
every save) computes $\Delta$ from consecutive snapshots and may never retire a
feature or amend prose --- a feature that looks emptied mid-edit usually
rebinds on the next save. The \emph{authority} pass (run at agent-session
close, on daemon startup, and from explicit sync) computes $\Delta$ state-based
against $\B$, is idempotent, recovers missed cycles, and is the only pass
allowed to raise \textsf{retire} and stale-prose \textsf{amend} proposals.

\subsection{Loop B: Realizing Tree Edits into Code}
\label{sec:loop-b}

Loop~B (Algorithm~\ref{alg:loop-b}) closes the opposite gap. Edits to the
rendered tree document are parsed and diffed against the store; verdicts on
pending proposals are drained from the review inbox; and live document-ahead
suggestions are applied by the loop itself (the human's only verb on their own
suggestion is \emph{withdraw} --- application is the agent side's
acknowledgment, stamped $\mu = \suggest$ with $\delta$ = the suggestion id).

Two ordering invariants keep text and store coherent
(lines~\ref{line:snapshot} and~\ref{line:rerender}). The text diff is
snapshotted \emph{before} any store mutation: verdict accepts move the store
ahead of the on-disk text, and diffing afterwards would misread the stale text
as a human edit that reverts the accepted change. Symmetrically, any pass that
mutated the store re-renders the text before returning, so the next pass diffs
a caught-up document.

Not every tree edit means ``write code.'' The \emph{imperative gate}
(line~\ref{line:imperative}) classifies each applied edit: descriptive prose
(``validates the token'') merely persists, while imperative prose (``should
validate the token'', ``Add retry logic''), an accepted unrealized plan node,
or retiring a feature that owns code each mint a \emph{directive} $d$ --- a
natural-language work order assembled from the feature's description and its
bound symbols, scoped to the files the feature owns. Directives are queued for
the user's \emph{live} coding-agent session rather than a headless model: the
session implements each directive, reflects the new code back through the
agent-side reflection API with $\delta = d$, and deletes the queue. Loop~A
then closes the cycle, with every resulting ledger entry carrying the causal
chain from the original document edit.

Verdict semantics are deliberately asymmetric for the destructive case
(line~\ref{line:retire}): accepting a machine-raised \textsf{retire} from the
inbox is \emph{detach-only} --- the feature is hidden and its chunks freed for
re-attribution, but no code-deletion directive is ever queued, because a retire
raised from transient drift may be a false positive and deleting code on a
single click is the most destructive failure mode. Only an explicit human
retire edit in the document text (or an agent proposal explicitly flagged to
delete code) queues removal work.

\begin{algorithm}[t]
\caption{Loop B --- tree $\to$ code realization (one pass)}
\label{alg:loop-b}
\begin{algorithmic}[1]
\Require tree text $T_{\mathrm{txt}}$; inbox verdicts $V$; intents $I$; store $(\T, \B, L)$
\Ensure updated $(\T, \B, L)$; directive queue $D$; refreshed $\sigma$
\State $\mathit{ops}_{\mathrm{user}} \gets \textsc{Diff}(T_{\mathrm{txt}}, \T)$ \Comment{snapshot BEFORE mutations (invariant 1)} \label{line:snapshot}
\State $D \gets \emptyset$
\For{$(e, \mathit{accept}) \in V$} \Comment{drain human verdicts on pending proposals}
    \If{$\mathit{accept}$}
        \State \textsc{Apply}$(e.\mathit{op}, \alpha{=}\textsf{human}, \mu{=}\pen)$; \textsc{Consume}$(e)$
        \If{$e.\mathit{op}$ is \textsf{retire} $\wedge\, \neg e.\mathit{op}.\mathit{deleteCode}$} \label{line:retire}
            \State \textsc{DetachAll}$(\B^{-1}(e.\mathit{op}.f))$ \Comment{detach-only: never queue code deletion}
        \ElsIf{\textsc{ImpliesCode}$(e.\mathit{op})$} $D \gets D \cup \{\textsc{Directive}(e.\mathit{op},\, \delta{=}e)\}$
        \EndIf
    \Else\ \textsc{Consume}$(e)$ \Comment{reject: delete the proposal}
    \EndIf
\EndFor
\For{$o \in \mathit{ops}_{\mathrm{user}}$} \Comment{direct edits are intentional: apply immediately}
    \State \textsc{Apply}$(o, \alpha, \mu \text{ from annotation channel (default \textsf{human}/\pen)})$
    \If{\textsc{ImpliesCode}$(o)$} $D \gets D \cup \{\textsc{Directive}(o, \delta{=}\ell_o)\}$ \label{line:imperative}
    \EndIf
\EndFor
\For{$i \in I$ with payload, fresh, unsatisfied} \Comment{doc-ahead suggestions: the loop applies them}
    \State \textsc{Apply}(\textsf{amend}$(i)$, $\alpha{=}i.\alpha$, $\mu{=}\suggest$, $\delta{=}i.\mathit{id}$)
    \If{\textsc{ImpliesCode}$(i)$} $D \gets D \cup \{\textsc{Directive}(i, \delta{=}i.\mathit{id})\}$
    \EndIf
\EndFor
\If{store mutated} \textsc{Render}$(\T) \to T_{\mathrm{txt}}$ \Comment{invariant 2: text catches up} \label{line:rerender}
\EndIf
\If{$D \neq \emptyset$}
    \State mint ids for $D$; queue $D$ for the live agent session; $H \gets H \cup \{d.f : d \in D\}$
    \State $\sigma \gets \awaiting$ \Comment{the agent implements, reflects with $\delta = d$, clears the queue}
\Else
    \State $\sigma \gets \textsc{RefreshState}()$
\EndIf
\State \Return $(\T, \B, L)$, $D$, $\sigma$
\end{algorithmic}
\end{algorithm}

\subsection{Conflict Resolution: the Ledger and Doc-Wins Holds}
\label{sec:conflicts}

Both loops route every decision through one classification table
(Table~\ref{tab:classify}): each detected change, from either side and by any
actor, maps to exactly one reaction. Three properties follow.

\paragraph{Provenance.} Every mutation is a ledger entry
(Definition~\ref{def:ledger}). The actor/mode pair separates \emph{who} changed
something from \emph{how committed} the change is, which drives both review
routing (only $\mu = \suggest$ items need verdicts) and the interface's
authorship rendering; $\delta$ threads causality across the realize cycle, so
a code change surfaced back by Loop~A is grouped under the document edit that
caused it.

\paragraph{Doc-wins holds.} Conflict resolution is a deterministic priority
rule rather than a merge heuristic. While $f \in H$ (a suggestion or queued
directive is pending on $f$), code-side \emph{intent} operations on $f$ ---
\textsf{amend}, \textsf{retire}, \textsf{move} --- are suppressed:
\begin{equation}
\label{eq:docwins}
\mathrm{suppress}(o, H) \;\iff\;
o.\mathit{kind} \in \{\textsf{amend}, \textsf{retire}, \textsf{move}\}
\;\wedge\; o.f \in H .
\end{equation}
Binding maintenance (\textsf{attach}/\textsf{detach}/\textsf{refresh}) is never
suppressed --- bindings are attribution, not intent, and must stay correct
regardless of who is mid-edit. Holds release when the suggestion is drained or
the directive queue completes, with a staleness backstop so an abandoned
suggestion cannot hold a feature indefinitely. The rule is the
tree-is-authoritative analogue of last-writer-wins: the side that expresses
\emph{intent} (the document) always outranks the side that merely
\emph{observes} (the code reflector).

\paragraph{Verdict grammar.} Who resolves what is structural:
accept/reject is exclusively a \emph{human} verb on \emph{code-ahead} items
(machine proposals), while a \emph{doc-ahead} suggestion is applied by the
machine side and offers its human author only \emph{withdraw}. No state is
double-resolvable, which eliminates accept/revert races by construction.

\begin{table}[t]
\centering
\caption{The change-classification table. Every change, from either side and by
any actor, maps to one reaction. $\auto$ reactions apply immediately;
$\suggest$ reactions await a human verdict; \textbf{D} marks reactions that
mint a realization directive.}
\label{tab:classify}
\small
\setlength{\tabcolsep}{4.5pt}
\begin{tabular}{@{}cllll@{}}
\toprule
\# & Side & Detected change & Origin & Reaction \\
\midrule
1 & code & modified bound chunk & any & \textsf{refresh} binding ($\auto$) \\
2 & code & removed bound chunk & any & \textsf{detach} ($\auto$); emptied $\wedge\, f \notin H \Rightarrow$ \textsf{retire} ($\suggest$, auth.\ pass) \\
3 & code & relocated chunk (Def.~\ref{def:reloc}) & any & deterministic \textsf{attach} to prior feature ($\auto$) \\
4 & code & added unbound chunk & any & LLM placement ($\auto$/$\suggest$); fallback $\mathrm{nbf}$ / \textsf{add} ($\suggest$) \\
5 & code & in-place modify, realized + prose & any & stale-prose \textsf{amend} (small $\to \auto$, large $\to \suggest$) \\
6 & code & any change during realize epoch & agent & rows 1--5 with $\delta = $ directive id \\
\midrule
7 & doc & descriptive amend / title edit & human, $\pen$ & apply immediately; no directive \\
8 & doc & imperative amend / plan add / retire-with-code & human, $\pen$ & apply + mint directive (\textbf{D}) \\
9 & doc & any edit & human, $\suggest$ & pending intent (hold); applied by Loop B drain $\to$ rows 7/8 \\
10 & doc & structural op via reflection API & agent & pending proposal ($\suggest$) \\
11 & doc & safe op via reflection API & agent & auto-apply; ledger $\alpha = \textsf{agent}$, $\mu = \auto$ \\
\midrule
12 & verdict & accept / reject & human & apply+consume / delete (\textsf{retire} accept is detach-only) \\
13 & code & drift on held feature ($f \in H$) & any & bindings maintained; \textsf{amend}/\textsf{retire}/\textsf{move} suppressed (Eq.~\ref{eq:docwins}) \\
\bottomrule
\end{tabular}
\end{table}

\subsection{Bootstrap}
\label{sec:bootstrap}

The initial tree is built by two scoped passes rather than one global
clustering. A \emph{per-file} pass issues one model call per source file,
seeing only that file's chunks and their incident graph edges --- making
cross-file ``junk drawer'' features structurally impossible --- and a
deterministic coverage step folds any chunk the model missed into the file's
largest feature. An \emph{organization} pass then groups the resulting
file-level features under a small number of thematic parents, using
feature-to-feature coupling aggregated from $\G$ as evidence. Both passes emit
ordinary ledger events through the same application seam as the loops, so
bootstrap output is curated (and corrected) with exactly the same verdict
machinery as any later change.
 from the ICLR template main file.
% Requires (all standard in the ICLR template ecosystem):
%   amsmath, amssymb, amsthm, booktabs, algorithm, algpseudocode
% plus the \newtheorem{definition}{Definition} declaration and the
% `problembox` environment + notation macros defined in main.tex.
% ------------------------------------------------------------------------------

\section{The \codoc{} System}
\label{sec:method}

\codoc{} maintains a human-authored \emph{feature tree} --- a hierarchy of named
units of intent --- bidirectionally synchronized with the code of a repository.
Unlike repository-level code graphs built \emph{from} code as retrieval context
for agents \citep{repograph2025,locagent2025,graphcoder2024,codexgraph2024}, the
tree is first-class authored intent: it is never regenerated from code, and the
attribution of code to features is a secondary, continuously repaired index.
Synchronization is asymmetric --- code changes are \emph{reflected} into the tree
(Loop~A, \S\ref{sec:loop-a}), while tree edits are \emph{realized} into code by a
coding agent (Loop~B, \S\ref{sec:loop-b}) --- and conflict resolution is
deterministic: pending document intent always wins over code-side observations
(\S\ref{sec:conflicts}).

\begin{problembox}{Problem statement}
Given a repository whose indexed chunk set is $\C$ and a human-authored feature
tree $\T = (\F, \prec)$, maintain (i)~a binding map
$\B \colon \C \to \F \cup \{\bot\}$ attributing every chunk to at most one
feature, and (ii)~the invariant that $\T$ expresses current human intent over
the code, under concurrent edits to either side by humans and autonomous
agents. Every mutation must be attributable (who, with what commitment, caused
by what), and divergence must resolve to a synchronization state
$\sigma = \synced$ without losing authored intent.
\end{problembox}

\subsection{Preliminaries}
\label{sec:prelim}

\begin{definition}[Chunk set]
\label{def:chunks}
The repository is indexed into a set of AST-level chunks
$\C = \{c_1, \dots, c_n\}$, where each chunk
$c = (\mathit{file},\, \mathit{sym},\, h^{\mathrm{tok}},\, h^{\mathrm{ast}})$
carries its file path, its qualified symbol path within the file, a
\emph{content hash} $h^{\mathrm{tok}}$ over the comment- and
whitespace-normalized token stream, and an \emph{AST-shape hash}
$h^{\mathrm{ast}}$ over the token-erased tree structure. The pair
$(\mathit{file}, \mathit{sym})$ is the chunk's identity key; the two hashes are
its relocation signals (Definition~\ref{def:reloc}).
\end{definition}

\begin{definition}[Feature tree]
\label{def:tree}
The feature tree is $\T = (\F, \prec)$, where $\F$ is a finite set of features
$f = (\mathit{id},\, \mathit{title},\, \mathit{desc},\, \mathit{realized})$ and
$\prec\, \subseteq \F \times \F$ is a forest-shaped parent relation. Titles and
descriptions are free prose authored by humans (or proposed by agents and
accepted by humans); $\mathit{realized} \in \{0,1\}$ marks whether the feature
is backed by code yet ($\mathit{realized}=0$ denotes an accepted plan
placeholder, flipped to $1$ by the first binding).
\end{definition}

\begin{definition}[Binding map]
\label{def:binding}
Attribution is a partial function $\B \colon \C \to \F \cup \{\bot\}$: each
chunk binds to at most one feature (enforced as a uniqueness constraint on the
chunk identity key), while one feature typically owns many chunks across many
files. Each binding stores the bound chunk's $h^{\mathrm{tok}}$ as its
\emph{fingerprint}, making staleness locally detectable. We write
$\B^{-1}(f) \subseteq \C$ for the chunks a feature owns and define coverage
$\mathrm{cov}(\B) = \lvert\{c : \B(c) \neq \bot\}\rvert / \lvert\C\rvert$.
\end{definition}

\begin{definition}[Code graph]
\label{def:graph}
Following the heterogeneous-graph convention of repository-level code
representations \citep{locagent2025,graphcoder2024}, the dependency structure
of the code is $\G = (\V, \E, \A, \R)$ with node-type map
$\tau \colon \V \to \A$, $\A = \{\textit{file}, \textit{class},
\textit{function}\}$, and edge-type map $\varphi \colon \E \to \R$,
$\R = \{\textit{contain}, \textit{call}, \textit{import}, \textit{inherit}\}$.
$\G$ is a derived cache, rebuildable from $\C$ at any time. The binding map
lifts $\G$ to a feature-coupling graph and defines the placement heuristic
\begin{equation}
\label{eq:neighbor}
\mathrm{nbf}(c) \;=\; \argmax_{f \in \F}\;
\bigl\lvert \{\, (v, v') \in \E \;:\; v' = \mathrm{sym}(c),\;
\B(\mathrm{chunk}(v)) = f \,\} \bigr\rvert ,
\end{equation}
the feature owning the most call/import edges incident to a new chunk $c$ ---
the deterministic fallback of Loop~A's coverage net
(Algorithm~\ref{alg:loop-a}, lines~\ref{line:covnet-start}--\ref{line:covnet-end}).
\end{definition}

\begin{definition}[Changeset]
\label{def:changeset}
A pass over the index yields $\Delta = (\Delta^{+}, \Delta^{-}, \Delta^{\sim})$:
chunks \emph{added} (present, unbound), \emph{removed} (a binding whose key has
left the index), and \emph{modified} (bound, but
$h^{\mathrm{tok}} \neq$ fingerprint). \codoc{} computes $\Delta$ in two modes:
a \emph{temporal} diff of consecutive index snapshots (cheap, runs on every
save) and a \emph{state-based} reconciliation that compares the index directly
against $\B$ (idempotent and self-healing; the authority pass). Only the latter
may retire features or amend prose, so destructive judgments are never made
from a transient mid-edit view.
\end{definition}

\begin{definition}[Relocation]
\label{def:reloc}
For $(r, a) \in \Delta^{-} \times \Delta^{+}$: $a$ is a \emph{move} of $r$ iff
$h^{\mathrm{tok}}_a = h^{\mathrm{tok}}_r$ (identical content at a new key), and
a \emph{rename} of $r$ iff $h^{\mathrm{ast}}_a = h^{\mathrm{ast}}_r$ uniquely
within the same file (same shape, new name; $1{:}1$ matches only). Relocations
re-attach attribution deterministically --- no model call, no risk of dropped
or duplicated attribution.
\end{definition}

\begin{definition}[Change ledger]
\label{def:ledger}
Every mutation of $(\T, \B)$ is an event in an append-only ledger $L$. A ledger
entry $\ell = (\mathit{op},\, f,\, \alpha,\, \mu,\, \delta)$ records the
operation, its target, the \emph{actor}
$\alpha \in \{\textsf{human}\} \cup \mathit{Agents} \cup \{\textsf{loop}\}$,
the \emph{mode} $\mu \in \{\pen, \suggest, \auto\}$ (committed edit, reviewable
proposal, or mechanical maintenance), and a \emph{causality link} $\delta$ ---
the directive, suggestion, or event id this change implements ($\bot$ if
spontaneous). A \emph{proposal} is an unapplied event; accepting it applies the
operation and consumes the event.
\end{definition}

\begin{definition}[Hold set and synchronization state]
\label{def:hold}
$H \subseteq \F$ is the set of features with pending document-ahead intent: a
live suggestion not yet drained, or a queued directive not yet implemented.
The synchronization state is
$\sigma \in \{\synced, \codedrift, \treedirty, \awaiting, \realizing\}$,
derived (not stored): pending code-side proposals put $\sigma = \codedrift$, a
non-empty directive queue forces $\sigma = \awaiting$ (a later code-side pass
cannot silently clobber queued work), and an active implementation pass reports
$\sigma = \realizing$.
\end{definition}

Table~\ref{tab:notation} summarizes the notation.

\begin{table}[t]
\centering
\caption{Notation.}
\label{tab:notation}
\small
\begin{tabular}{@{}ll@{}}
\toprule
Symbol & Meaning \\
\midrule
$\C$ & chunk set; $c = (\mathit{file}, \mathit{sym}, h^{\mathrm{tok}}, h^{\mathrm{ast}})$ \\
$\T = (\F, \prec)$ & feature tree (features + forest parent relation) \\
$\B \colon \C \to \F \cup \{\bot\}$ & binding map (each chunk owned by $\leq 1$ feature) \\
$\G = (\V, \E, \A, \R)$ & heterogeneous code dependency graph (derived cache) \\
$\Delta = (\Delta^{+}, \Delta^{-}, \Delta^{\sim})$ & changeset: added / removed / modified chunks \\
$h^{\mathrm{tok}}, h^{\mathrm{ast}}$ & content hash (move-invariant), AST-shape hash (rename-invariant) \\
$L$; $\ell = (\mathit{op}, f, \alpha, \mu, \delta)$ & append-only ledger; entry = op, target, actor, mode, cause \\
$H \subseteq \F$ & hold set (features with pending doc-ahead intent) \\
$\sigma$ & synchronization state \\
$\mathrm{nbf}(c)$ & graph-neighbor feature of a chunk (Eq.~\ref{eq:neighbor}) \\
$D$; $d$ & directive queue; a single realization directive \\
\bottomrule
\end{tabular}
\end{table}

\subsection{Loop A: Reflecting Code into the Tree}
\label{sec:loop-a}

Loop~A (Algorithm~\ref{alg:loop-a}) closes the gap from code to intent. Its
design principle is \emph{determinism first}: every change that can be resolved
mechanically --- refreshing fingerprints, detaching deleted chunks, carrying
attribution across moves and renames --- is resolved without a model
(lines~\ref{line:auto-start}--\ref{line:reloc-end}). A single LLM pass
(line~\ref{line:llm}) handles only the judgment calls that remain: placing
genuinely new code, repairing features that lost their last binding, and (in
the authority pass) refreshing prose whose code changed underneath it. The
pass sees the residual changeset, a file-locality subtree of $\T$, and
\emph{every} feature title; supplying the full title set to one call is what
prevents duplicate features without any post-hoc deduplication machinery.

Model output is not trusted with destructive authority. Proposed operations
are split by safety: binding maintenance and small description amendments
(bounded by an edit-distance ratio) auto-apply with $\mu = \auto$; structural
operations (add / move / retire) become unapplied proposal events with
$\mu = \suggest$ that a human accepts or rejects. Operations targeting held
features are dropped (line~\ref{line:hold}; \S\ref{sec:conflicts}), and a
\emph{coverage net}
(lines~\ref{line:covnet-start}--\ref{line:covnet-end}) guarantees no added
chunk is silently dropped: anything the model failed to place attaches to its
graph-neighbor feature $\mathrm{nbf}(c)$ or surfaces as an explicit
\textsf{add} proposal. Finally, stale-proposal garbage collection keeps the
pass idempotent: a pending \textsf{add} whose chunks have since bound
elsewhere, or a pending \textsf{retire} whose feature regained code, is
dropped, so repeated passes converge to $\sigma = \synced$ instead of
accumulating false positives.

\begin{algorithm}[t]
\caption{Loop A --- code $\to$ tree reflection (one pass)}
\label{alg:loop-a}
\begin{algorithmic}[1]
\Require index snapshot of $\C$; store $(\T, \B, L)$; hold set $H$; authority flag $\mathit{auth}$
\Ensure updated $(\T, \B, L)$; refreshed $\sigma$
\State $\Delta \gets \textsc{Changeset}(\C, \B)$ \Comment{temporal diff, or state-based iff $\mathit{auth}$}
\State \textsc{GcStaleProposals}$(L, \Delta)$ \Comment{drop superseded \textsf{add}/\textsf{retire} proposals}
\For{$c \in \Delta^{\sim}$} \textsc{Apply}(\textsf{refresh}$(c)$, $\mu{=}\auto$) \EndFor \label{line:auto-start}
\For{$c \in \Delta^{-}$} \textsc{Apply}(\textsf{detach}$(c)$, $\mu{=}\auto$) \EndFor
\For{$(r, a) \in \textsc{Relocations}(\Delta^{-}, \Delta^{+})$} \Comment{Def.~\ref{def:reloc}: $h^{\mathrm{tok}}$ move / $h^{\mathrm{ast}}$ rename}
    \State \textsc{Apply}(\textsf{attach}$(a, \B(r))$, $\mu{=}\auto$) \label{line:reloc-end}
\EndFor
\State $U \gets \{c \in \Delta^{+} : \B(c) = \bot\}$;\quad
       $E \gets \{f : \B^{-1}(f) = \emptyset \wedge f \notin H\}$ \Comment{unbound adds; emptied features}
\If{$U \neq \emptyset \,\vee\, E \neq \emptyset \,\vee\, (\mathit{auth} \wedge \textsc{StaleProse}(\Delta^{\sim}, H))$}
    \State $\mathit{ops} \gets \textsc{LLM}\bigl(\Delta,\ \textsc{Subtree}(\T, \Delta),\ \textsc{Titles}(\T),\ \G\bigr)$ \Comment{one call per pass} \label{line:llm}
    \State $\mathit{ops} \gets \{o \in \mathit{ops} : \neg\,\textsc{SuppressedByHold}(o, H)\}$ \label{line:hold}
    \For{$o \in \mathit{ops}$}
        \If{\textsc{Safe}$(o)$} \textsc{Apply}$(o, \mu{=}\auto)$
        \Else\ \textsc{Propose}$(o, \mu{=}\suggest)$ \Comment{unapplied event; human verdict pending}
        \EndIf
    \EndFor
\EndIf
\For{$c \in U$ with $\B(c) = \bot$ still} \label{line:covnet-start} \Comment{coverage net: nothing is dropped}
    \If{$\mathrm{nbf}(c) \neq \bot$} \textsc{Apply}(\textsf{attach}$(c, \mathrm{nbf}(c))$, $\mu{=}\auto$)
    \Else\ \textsc{Propose}(\textsf{add}$(c)$, $\mu{=}\suggest$)
    \EndIf
\EndFor \label{line:covnet-end}
\State \Return $(\T, \B, L)$, \textsc{RefreshState}$()$
\end{algorithmic}
\end{algorithm}

Two authority levels share this pipeline. The \emph{temporal} pass (run on
every save) computes $\Delta$ from consecutive snapshots and may never retire a
feature or amend prose --- a feature that looks emptied mid-edit usually
rebinds on the next save. The \emph{authority} pass (run at agent-session
close, on daemon startup, and from explicit sync) computes $\Delta$ state-based
against $\B$, is idempotent, recovers missed cycles, and is the only pass
allowed to raise \textsf{retire} and stale-prose \textsf{amend} proposals.

\subsection{Loop B: Realizing Tree Edits into Code}
\label{sec:loop-b}

Loop~B (Algorithm~\ref{alg:loop-b}) closes the opposite gap. Edits to the
rendered tree document are parsed and diffed against the store; verdicts on
pending proposals are drained from the review inbox; and live document-ahead
suggestions are applied by the loop itself (the human's only verb on their own
suggestion is \emph{withdraw} --- application is the agent side's
acknowledgment, stamped $\mu = \suggest$ with $\delta$ = the suggestion id).

Two ordering invariants keep text and store coherent
(lines~\ref{line:snapshot} and~\ref{line:rerender}). The text diff is
snapshotted \emph{before} any store mutation: verdict accepts move the store
ahead of the on-disk text, and diffing afterwards would misread the stale text
as a human edit that reverts the accepted change. Symmetrically, any pass that
mutated the store re-renders the text before returning, so the next pass diffs
a caught-up document.

Not every tree edit means ``write code.'' The \emph{imperative gate}
(line~\ref{line:imperative}) classifies each applied edit: descriptive prose
(``validates the token'') merely persists, while imperative prose (``should
validate the token'', ``Add retry logic''), an accepted unrealized plan node,
or retiring a feature that owns code each mint a \emph{directive} $d$ --- a
natural-language work order assembled from the feature's description and its
bound symbols, scoped to the files the feature owns. Directives are queued for
the user's \emph{live} coding-agent session rather than a headless model: the
session implements each directive, reflects the new code back through the
agent-side reflection API with $\delta = d$, and deletes the queue. Loop~A
then closes the cycle, with every resulting ledger entry carrying the causal
chain from the original document edit.

Verdict semantics are deliberately asymmetric for the destructive case
(line~\ref{line:retire}): accepting a machine-raised \textsf{retire} from the
inbox is \emph{detach-only} --- the feature is hidden and its chunks freed for
re-attribution, but no code-deletion directive is ever queued, because a retire
raised from transient drift may be a false positive and deleting code on a
single click is the most destructive failure mode. Only an explicit human
retire edit in the document text (or an agent proposal explicitly flagged to
delete code) queues removal work.

\begin{algorithm}[t]
\caption{Loop B --- tree $\to$ code realization (one pass)}
\label{alg:loop-b}
\begin{algorithmic}[1]
\Require tree text $T_{\mathrm{txt}}$; inbox verdicts $V$; intents $I$; store $(\T, \B, L)$
\Ensure updated $(\T, \B, L)$; directive queue $D$; refreshed $\sigma$
\State $\mathit{ops}_{\mathrm{user}} \gets \textsc{Diff}(T_{\mathrm{txt}}, \T)$ \Comment{snapshot BEFORE mutations (invariant 1)} \label{line:snapshot}
\State $D \gets \emptyset$
\For{$(e, \mathit{accept}) \in V$} \Comment{drain human verdicts on pending proposals}
    \If{$\mathit{accept}$}
        \State \textsc{Apply}$(e.\mathit{op}, \alpha{=}\textsf{human}, \mu{=}\pen)$; \textsc{Consume}$(e)$
        \If{$e.\mathit{op}$ is \textsf{retire} $\wedge\, \neg e.\mathit{op}.\mathit{deleteCode}$} \label{line:retire}
            \State \textsc{DetachAll}$(\B^{-1}(e.\mathit{op}.f))$ \Comment{detach-only: never queue code deletion}
        \ElsIf{\textsc{ImpliesCode}$(e.\mathit{op})$} $D \gets D \cup \{\textsc{Directive}(e.\mathit{op},\, \delta{=}e)\}$
        \EndIf
    \Else\ \textsc{Consume}$(e)$ \Comment{reject: delete the proposal}
    \EndIf
\EndFor
\For{$o \in \mathit{ops}_{\mathrm{user}}$} \Comment{direct edits are intentional: apply immediately}
    \State \textsc{Apply}$(o, \alpha, \mu \text{ from annotation channel (default \textsf{human}/\pen)})$
    \If{\textsc{ImpliesCode}$(o)$} $D \gets D \cup \{\textsc{Directive}(o, \delta{=}\ell_o)\}$ \label{line:imperative}
    \EndIf
\EndFor
\For{$i \in I$ with payload, fresh, unsatisfied} \Comment{doc-ahead suggestions: the loop applies them}
    \State \textsc{Apply}(\textsf{amend}$(i)$, $\alpha{=}i.\alpha$, $\mu{=}\suggest$, $\delta{=}i.\mathit{id}$)
    \If{\textsc{ImpliesCode}$(i)$} $D \gets D \cup \{\textsc{Directive}(i, \delta{=}i.\mathit{id})\}$
    \EndIf
\EndFor
\If{store mutated} \textsc{Render}$(\T) \to T_{\mathrm{txt}}$ \Comment{invariant 2: text catches up} \label{line:rerender}
\EndIf
\If{$D \neq \emptyset$}
    \State mint ids for $D$; queue $D$ for the live agent session; $H \gets H \cup \{d.f : d \in D\}$
    \State $\sigma \gets \awaiting$ \Comment{the agent implements, reflects with $\delta = d$, clears the queue}
\Else
    \State $\sigma \gets \textsc{RefreshState}()$
\EndIf
\State \Return $(\T, \B, L)$, $D$, $\sigma$
\end{algorithmic}
\end{algorithm}

\subsection{Conflict Resolution: the Ledger and Doc-Wins Holds}
\label{sec:conflicts}

Both loops route every decision through one classification table
(Table~\ref{tab:classify}): each detected change, from either side and by any
actor, maps to exactly one reaction. Three properties follow.

\paragraph{Provenance.} Every mutation is a ledger entry
(Definition~\ref{def:ledger}). The actor/mode pair separates \emph{who} changed
something from \emph{how committed} the change is, which drives both review
routing (only $\mu = \suggest$ items need verdicts) and the interface's
authorship rendering; $\delta$ threads causality across the realize cycle, so
a code change surfaced back by Loop~A is grouped under the document edit that
caused it.

\paragraph{Doc-wins holds.} Conflict resolution is a deterministic priority
rule rather than a merge heuristic. While $f \in H$ (a suggestion or queued
directive is pending on $f$), code-side \emph{intent} operations on $f$ ---
\textsf{amend}, \textsf{retire}, \textsf{move} --- are suppressed:
\begin{equation}
\label{eq:docwins}
\mathrm{suppress}(o, H) \;\iff\;
o.\mathit{kind} \in \{\textsf{amend}, \textsf{retire}, \textsf{move}\}
\;\wedge\; o.f \in H .
\end{equation}
Binding maintenance (\textsf{attach}/\textsf{detach}/\textsf{refresh}) is never
suppressed --- bindings are attribution, not intent, and must stay correct
regardless of who is mid-edit. Holds release when the suggestion is drained or
the directive queue completes, with a staleness backstop so an abandoned
suggestion cannot hold a feature indefinitely. The rule is the
tree-is-authoritative analogue of last-writer-wins: the side that expresses
\emph{intent} (the document) always outranks the side that merely
\emph{observes} (the code reflector).

\paragraph{Verdict grammar.} Who resolves what is structural:
accept/reject is exclusively a \emph{human} verb on \emph{code-ahead} items
(machine proposals), while a \emph{doc-ahead} suggestion is applied by the
machine side and offers its human author only \emph{withdraw}. No state is
double-resolvable, which eliminates accept/revert races by construction.

\begin{table}[t]
\centering
\caption{The change-classification table. Every change, from either side and by
any actor, maps to one reaction. $\auto$ reactions apply immediately;
$\suggest$ reactions await a human verdict; \textbf{D} marks reactions that
mint a realization directive.}
\label{tab:classify}
\small
\setlength{\tabcolsep}{4.5pt}
\begin{tabular}{@{}cllll@{}}
\toprule
\# & Side & Detected change & Origin & Reaction \\
\midrule
1 & code & modified bound chunk & any & \textsf{refresh} binding ($\auto$) \\
2 & code & removed bound chunk & any & \textsf{detach} ($\auto$); emptied $\wedge\, f \notin H \Rightarrow$ \textsf{retire} ($\suggest$, auth.\ pass) \\
3 & code & relocated chunk (Def.~\ref{def:reloc}) & any & deterministic \textsf{attach} to prior feature ($\auto$) \\
4 & code & added unbound chunk & any & LLM placement ($\auto$/$\suggest$); fallback $\mathrm{nbf}$ / \textsf{add} ($\suggest$) \\
5 & code & in-place modify, realized + prose & any & stale-prose \textsf{amend} (small $\to \auto$, large $\to \suggest$) \\
6 & code & any change during realize epoch & agent & rows 1--5 with $\delta = $ directive id \\
\midrule
7 & doc & descriptive amend / title edit & human, $\pen$ & apply immediately; no directive \\
8 & doc & imperative amend / plan add / retire-with-code & human, $\pen$ & apply + mint directive (\textbf{D}) \\
9 & doc & any edit & human, $\suggest$ & pending intent (hold); applied by Loop B drain $\to$ rows 7/8 \\
10 & doc & structural op via reflection API & agent & pending proposal ($\suggest$) \\
11 & doc & safe op via reflection API & agent & auto-apply; ledger $\alpha = \textsf{agent}$, $\mu = \auto$ \\
\midrule
12 & verdict & accept / reject & human & apply+consume / delete (\textsf{retire} accept is detach-only) \\
13 & code & drift on held feature ($f \in H$) & any & bindings maintained; \textsf{amend}/\textsf{retire}/\textsf{move} suppressed (Eq.~\ref{eq:docwins}) \\
\bottomrule
\end{tabular}
\end{table}

\subsection{Bootstrap}
\label{sec:bootstrap}

The initial tree is built by two scoped passes rather than one global
clustering. A \emph{per-file} pass issues one model call per source file,
seeing only that file's chunks and their incident graph edges --- making
cross-file ``junk drawer'' features structurally impossible --- and a
deterministic coverage step folds any chunk the model missed into the file's
largest feature. An \emph{organization} pass then groups the resulting
file-level features under a small number of thematic parents, using
feature-to-feature coupling aggregated from $\G$ as evidence. Both passes emit
ordinary ledger events through the same application seam as the loops, so
bootstrap output is curated (and corrected) with exactly the same verdict
machinery as any later change.
 where the method section goes;
% everything here is template-compatible (algorithmicx, amsthm, booktabs,
% tcolorbox are all available alongside iclr20xx.sty).
%
% Build:  pdflatex main && bibtex main && pdflatex main && pdflatex main
% ------------------------------------------------------------------------------
\documentclass[11pt]{article}

\usepackage[margin=1in]{geometry}
\usepackage{times}
\usepackage[hidelinks]{hyperref}
\usepackage[numbers]{natbib}

% ── [ICLR] packages required by method.tex ────────────────────────────────────
\usepackage{amsmath, amssymb, amsthm}
\usepackage{booktabs}
\usepackage{algorithm}
\usepackage{algpseudocode}

% Problem-statement box. Plain-LaTeX implementation (no tcolorbox dependency);
% if your ICLR build has tcolorbox, feel free to swap in a tcolorbox variant.
\newenvironment{problembox}[1]{%
  \par\medskip\noindent
  \begin{center}\begin{minipage}{0.96\linewidth}%
  \hrule height 0.8pt \vspace{4pt}%
  \noindent\textbf{#1.}\enspace\itshape
}{%
  \par\vspace{4pt}\hrule height 0.8pt
  \end{minipage}\end{center}\medskip
}

% ── [ICLR] theorem environments ───────────────────────────────────────────────
\theoremstyle{definition}
\newtheorem{definition}{Definition}

% ── [ICLR] macros used by method.tex ──────────────────────────────────────────
\newcommand{\codoc}{\textsc{codoc}}
\newcommand{\C}{\mathcal{C}}
\newcommand{\F}{\mathcal{F}}
\newcommand{\T}{\mathcal{T}}
\newcommand{\B}{\mathcal{B}}
\newcommand{\G}{\mathcal{G}}
\newcommand{\V}{\mathcal{V}}
\newcommand{\E}{\mathcal{E}}
\newcommand{\A}{\mathcal{A}}
\newcommand{\R}{\mathcal{R}}
\DeclareMathOperator*{\argmax}{arg\,max}
% mode / state names
\newcommand{\pen}{\textsf{pen}}
\newcommand{\suggest}{\textsf{suggest}}
\newcommand{\auto}{\textsf{auto}}
\newcommand{\synced}{\textsf{in\_sync}}
\newcommand{\codedrift}{\textsf{code\_drift}}
\newcommand{\treedirty}{\textsf{tree\_dirty}}
\newcommand{\awaiting}{\textsf{awaiting\_impl}}
\newcommand{\realizing}{\textsf{realizing}}

\title{\codoc{}: Bidirectional Synchronization of a Human-Intent Feature Tree
with Code\\ \large (method section draft)}
\author{Anonymous}
\date{}

\begin{document}
\maketitle

% ------------------------------------------------------------------------------
% method.tex — the codoc method/algorithm section.
%
% Self-contained: % ------------------------------------------------------------------------------
% method.tex — the codoc method/algorithm section.
%
% Self-contained: % ------------------------------------------------------------------------------
% method.tex — the codoc method/algorithm section.
%
% Self-contained: \input{method} from the ICLR template main file.
% Requires (all standard in the ICLR template ecosystem):
%   amsmath, amssymb, amsthm, booktabs, algorithm, algpseudocode
% plus the \newtheorem{definition}{Definition} declaration and the
% `problembox` environment + notation macros defined in main.tex.
% ------------------------------------------------------------------------------

\section{The \codoc{} System}
\label{sec:method}

\codoc{} maintains a human-authored \emph{feature tree} --- a hierarchy of named
units of intent --- bidirectionally synchronized with the code of a repository.
Unlike repository-level code graphs built \emph{from} code as retrieval context
for agents \citep{repograph2025,locagent2025,graphcoder2024,codexgraph2024}, the
tree is first-class authored intent: it is never regenerated from code, and the
attribution of code to features is a secondary, continuously repaired index.
Synchronization is asymmetric --- code changes are \emph{reflected} into the tree
(Loop~A, \S\ref{sec:loop-a}), while tree edits are \emph{realized} into code by a
coding agent (Loop~B, \S\ref{sec:loop-b}) --- and conflict resolution is
deterministic: pending document intent always wins over code-side observations
(\S\ref{sec:conflicts}).

\begin{problembox}{Problem statement}
Given a repository whose indexed chunk set is $\C$ and a human-authored feature
tree $\T = (\F, \prec)$, maintain (i)~a binding map
$\B \colon \C \to \F \cup \{\bot\}$ attributing every chunk to at most one
feature, and (ii)~the invariant that $\T$ expresses current human intent over
the code, under concurrent edits to either side by humans and autonomous
agents. Every mutation must be attributable (who, with what commitment, caused
by what), and divergence must resolve to a synchronization state
$\sigma = \synced$ without losing authored intent.
\end{problembox}

\subsection{Preliminaries}
\label{sec:prelim}

\begin{definition}[Chunk set]
\label{def:chunks}
The repository is indexed into a set of AST-level chunks
$\C = \{c_1, \dots, c_n\}$, where each chunk
$c = (\mathit{file},\, \mathit{sym},\, h^{\mathrm{tok}},\, h^{\mathrm{ast}})$
carries its file path, its qualified symbol path within the file, a
\emph{content hash} $h^{\mathrm{tok}}$ over the comment- and
whitespace-normalized token stream, and an \emph{AST-shape hash}
$h^{\mathrm{ast}}$ over the token-erased tree structure. The pair
$(\mathit{file}, \mathit{sym})$ is the chunk's identity key; the two hashes are
its relocation signals (Definition~\ref{def:reloc}).
\end{definition}

\begin{definition}[Feature tree]
\label{def:tree}
The feature tree is $\T = (\F, \prec)$, where $\F$ is a finite set of features
$f = (\mathit{id},\, \mathit{title},\, \mathit{desc},\, \mathit{realized})$ and
$\prec\, \subseteq \F \times \F$ is a forest-shaped parent relation. Titles and
descriptions are free prose authored by humans (or proposed by agents and
accepted by humans); $\mathit{realized} \in \{0,1\}$ marks whether the feature
is backed by code yet ($\mathit{realized}=0$ denotes an accepted plan
placeholder, flipped to $1$ by the first binding).
\end{definition}

\begin{definition}[Binding map]
\label{def:binding}
Attribution is a partial function $\B \colon \C \to \F \cup \{\bot\}$: each
chunk binds to at most one feature (enforced as a uniqueness constraint on the
chunk identity key), while one feature typically owns many chunks across many
files. Each binding stores the bound chunk's $h^{\mathrm{tok}}$ as its
\emph{fingerprint}, making staleness locally detectable. We write
$\B^{-1}(f) \subseteq \C$ for the chunks a feature owns and define coverage
$\mathrm{cov}(\B) = \lvert\{c : \B(c) \neq \bot\}\rvert / \lvert\C\rvert$.
\end{definition}

\begin{definition}[Code graph]
\label{def:graph}
Following the heterogeneous-graph convention of repository-level code
representations \citep{locagent2025,graphcoder2024}, the dependency structure
of the code is $\G = (\V, \E, \A, \R)$ with node-type map
$\tau \colon \V \to \A$, $\A = \{\textit{file}, \textit{class},
\textit{function}\}$, and edge-type map $\varphi \colon \E \to \R$,
$\R = \{\textit{contain}, \textit{call}, \textit{import}, \textit{inherit}\}$.
$\G$ is a derived cache, rebuildable from $\C$ at any time. The binding map
lifts $\G$ to a feature-coupling graph and defines the placement heuristic
\begin{equation}
\label{eq:neighbor}
\mathrm{nbf}(c) \;=\; \argmax_{f \in \F}\;
\bigl\lvert \{\, (v, v') \in \E \;:\; v' = \mathrm{sym}(c),\;
\B(\mathrm{chunk}(v)) = f \,\} \bigr\rvert ,
\end{equation}
the feature owning the most call/import edges incident to a new chunk $c$ ---
the deterministic fallback of Loop~A's coverage net
(Algorithm~\ref{alg:loop-a}, lines~\ref{line:covnet-start}--\ref{line:covnet-end}).
\end{definition}

\begin{definition}[Changeset]
\label{def:changeset}
A pass over the index yields $\Delta = (\Delta^{+}, \Delta^{-}, \Delta^{\sim})$:
chunks \emph{added} (present, unbound), \emph{removed} (a binding whose key has
left the index), and \emph{modified} (bound, but
$h^{\mathrm{tok}} \neq$ fingerprint). \codoc{} computes $\Delta$ in two modes:
a \emph{temporal} diff of consecutive index snapshots (cheap, runs on every
save) and a \emph{state-based} reconciliation that compares the index directly
against $\B$ (idempotent and self-healing; the authority pass). Only the latter
may retire features or amend prose, so destructive judgments are never made
from a transient mid-edit view.
\end{definition}

\begin{definition}[Relocation]
\label{def:reloc}
For $(r, a) \in \Delta^{-} \times \Delta^{+}$: $a$ is a \emph{move} of $r$ iff
$h^{\mathrm{tok}}_a = h^{\mathrm{tok}}_r$ (identical content at a new key), and
a \emph{rename} of $r$ iff $h^{\mathrm{ast}}_a = h^{\mathrm{ast}}_r$ uniquely
within the same file (same shape, new name; $1{:}1$ matches only). Relocations
re-attach attribution deterministically --- no model call, no risk of dropped
or duplicated attribution.
\end{definition}

\begin{definition}[Change ledger]
\label{def:ledger}
Every mutation of $(\T, \B)$ is an event in an append-only ledger $L$. A ledger
entry $\ell = (\mathit{op},\, f,\, \alpha,\, \mu,\, \delta)$ records the
operation, its target, the \emph{actor}
$\alpha \in \{\textsf{human}\} \cup \mathit{Agents} \cup \{\textsf{loop}\}$,
the \emph{mode} $\mu \in \{\pen, \suggest, \auto\}$ (committed edit, reviewable
proposal, or mechanical maintenance), and a \emph{causality link} $\delta$ ---
the directive, suggestion, or event id this change implements ($\bot$ if
spontaneous). A \emph{proposal} is an unapplied event; accepting it applies the
operation and consumes the event.
\end{definition}

\begin{definition}[Hold set and synchronization state]
\label{def:hold}
$H \subseteq \F$ is the set of features with pending document-ahead intent: a
live suggestion not yet drained, or a queued directive not yet implemented.
The synchronization state is
$\sigma \in \{\synced, \codedrift, \treedirty, \awaiting, \realizing\}$,
derived (not stored): pending code-side proposals put $\sigma = \codedrift$, a
non-empty directive queue forces $\sigma = \awaiting$ (a later code-side pass
cannot silently clobber queued work), and an active implementation pass reports
$\sigma = \realizing$.
\end{definition}

Table~\ref{tab:notation} summarizes the notation.

\begin{table}[t]
\centering
\caption{Notation.}
\label{tab:notation}
\small
\begin{tabular}{@{}ll@{}}
\toprule
Symbol & Meaning \\
\midrule
$\C$ & chunk set; $c = (\mathit{file}, \mathit{sym}, h^{\mathrm{tok}}, h^{\mathrm{ast}})$ \\
$\T = (\F, \prec)$ & feature tree (features + forest parent relation) \\
$\B \colon \C \to \F \cup \{\bot\}$ & binding map (each chunk owned by $\leq 1$ feature) \\
$\G = (\V, \E, \A, \R)$ & heterogeneous code dependency graph (derived cache) \\
$\Delta = (\Delta^{+}, \Delta^{-}, \Delta^{\sim})$ & changeset: added / removed / modified chunks \\
$h^{\mathrm{tok}}, h^{\mathrm{ast}}$ & content hash (move-invariant), AST-shape hash (rename-invariant) \\
$L$; $\ell = (\mathit{op}, f, \alpha, \mu, \delta)$ & append-only ledger; entry = op, target, actor, mode, cause \\
$H \subseteq \F$ & hold set (features with pending doc-ahead intent) \\
$\sigma$ & synchronization state \\
$\mathrm{nbf}(c)$ & graph-neighbor feature of a chunk (Eq.~\ref{eq:neighbor}) \\
$D$; $d$ & directive queue; a single realization directive \\
\bottomrule
\end{tabular}
\end{table}

\subsection{Loop A: Reflecting Code into the Tree}
\label{sec:loop-a}

Loop~A (Algorithm~\ref{alg:loop-a}) closes the gap from code to intent. Its
design principle is \emph{determinism first}: every change that can be resolved
mechanically --- refreshing fingerprints, detaching deleted chunks, carrying
attribution across moves and renames --- is resolved without a model
(lines~\ref{line:auto-start}--\ref{line:reloc-end}). A single LLM pass
(line~\ref{line:llm}) handles only the judgment calls that remain: placing
genuinely new code, repairing features that lost their last binding, and (in
the authority pass) refreshing prose whose code changed underneath it. The
pass sees the residual changeset, a file-locality subtree of $\T$, and
\emph{every} feature title; supplying the full title set to one call is what
prevents duplicate features without any post-hoc deduplication machinery.

Model output is not trusted with destructive authority. Proposed operations
are split by safety: binding maintenance and small description amendments
(bounded by an edit-distance ratio) auto-apply with $\mu = \auto$; structural
operations (add / move / retire) become unapplied proposal events with
$\mu = \suggest$ that a human accepts or rejects. Operations targeting held
features are dropped (line~\ref{line:hold}; \S\ref{sec:conflicts}), and a
\emph{coverage net}
(lines~\ref{line:covnet-start}--\ref{line:covnet-end}) guarantees no added
chunk is silently dropped: anything the model failed to place attaches to its
graph-neighbor feature $\mathrm{nbf}(c)$ or surfaces as an explicit
\textsf{add} proposal. Finally, stale-proposal garbage collection keeps the
pass idempotent: a pending \textsf{add} whose chunks have since bound
elsewhere, or a pending \textsf{retire} whose feature regained code, is
dropped, so repeated passes converge to $\sigma = \synced$ instead of
accumulating false positives.

\begin{algorithm}[t]
\caption{Loop A --- code $\to$ tree reflection (one pass)}
\label{alg:loop-a}
\begin{algorithmic}[1]
\Require index snapshot of $\C$; store $(\T, \B, L)$; hold set $H$; authority flag $\mathit{auth}$
\Ensure updated $(\T, \B, L)$; refreshed $\sigma$
\State $\Delta \gets \textsc{Changeset}(\C, \B)$ \Comment{temporal diff, or state-based iff $\mathit{auth}$}
\State \textsc{GcStaleProposals}$(L, \Delta)$ \Comment{drop superseded \textsf{add}/\textsf{retire} proposals}
\For{$c \in \Delta^{\sim}$} \textsc{Apply}(\textsf{refresh}$(c)$, $\mu{=}\auto$) \EndFor \label{line:auto-start}
\For{$c \in \Delta^{-}$} \textsc{Apply}(\textsf{detach}$(c)$, $\mu{=}\auto$) \EndFor
\For{$(r, a) \in \textsc{Relocations}(\Delta^{-}, \Delta^{+})$} \Comment{Def.~\ref{def:reloc}: $h^{\mathrm{tok}}$ move / $h^{\mathrm{ast}}$ rename}
    \State \textsc{Apply}(\textsf{attach}$(a, \B(r))$, $\mu{=}\auto$) \label{line:reloc-end}
\EndFor
\State $U \gets \{c \in \Delta^{+} : \B(c) = \bot\}$;\quad
       $E \gets \{f : \B^{-1}(f) = \emptyset \wedge f \notin H\}$ \Comment{unbound adds; emptied features}
\If{$U \neq \emptyset \,\vee\, E \neq \emptyset \,\vee\, (\mathit{auth} \wedge \textsc{StaleProse}(\Delta^{\sim}, H))$}
    \State $\mathit{ops} \gets \textsc{LLM}\bigl(\Delta,\ \textsc{Subtree}(\T, \Delta),\ \textsc{Titles}(\T),\ \G\bigr)$ \Comment{one call per pass} \label{line:llm}
    \State $\mathit{ops} \gets \{o \in \mathit{ops} : \neg\,\textsc{SuppressedByHold}(o, H)\}$ \label{line:hold}
    \For{$o \in \mathit{ops}$}
        \If{\textsc{Safe}$(o)$} \textsc{Apply}$(o, \mu{=}\auto)$
        \Else\ \textsc{Propose}$(o, \mu{=}\suggest)$ \Comment{unapplied event; human verdict pending}
        \EndIf
    \EndFor
\EndIf
\For{$c \in U$ with $\B(c) = \bot$ still} \label{line:covnet-start} \Comment{coverage net: nothing is dropped}
    \If{$\mathrm{nbf}(c) \neq \bot$} \textsc{Apply}(\textsf{attach}$(c, \mathrm{nbf}(c))$, $\mu{=}\auto$)
    \Else\ \textsc{Propose}(\textsf{add}$(c)$, $\mu{=}\suggest$)
    \EndIf
\EndFor \label{line:covnet-end}
\State \Return $(\T, \B, L)$, \textsc{RefreshState}$()$
\end{algorithmic}
\end{algorithm}

Two authority levels share this pipeline. The \emph{temporal} pass (run on
every save) computes $\Delta$ from consecutive snapshots and may never retire a
feature or amend prose --- a feature that looks emptied mid-edit usually
rebinds on the next save. The \emph{authority} pass (run at agent-session
close, on daemon startup, and from explicit sync) computes $\Delta$ state-based
against $\B$, is idempotent, recovers missed cycles, and is the only pass
allowed to raise \textsf{retire} and stale-prose \textsf{amend} proposals.

\subsection{Loop B: Realizing Tree Edits into Code}
\label{sec:loop-b}

Loop~B (Algorithm~\ref{alg:loop-b}) closes the opposite gap. Edits to the
rendered tree document are parsed and diffed against the store; verdicts on
pending proposals are drained from the review inbox; and live document-ahead
suggestions are applied by the loop itself (the human's only verb on their own
suggestion is \emph{withdraw} --- application is the agent side's
acknowledgment, stamped $\mu = \suggest$ with $\delta$ = the suggestion id).

Two ordering invariants keep text and store coherent
(lines~\ref{line:snapshot} and~\ref{line:rerender}). The text diff is
snapshotted \emph{before} any store mutation: verdict accepts move the store
ahead of the on-disk text, and diffing afterwards would misread the stale text
as a human edit that reverts the accepted change. Symmetrically, any pass that
mutated the store re-renders the text before returning, so the next pass diffs
a caught-up document.

Not every tree edit means ``write code.'' The \emph{imperative gate}
(line~\ref{line:imperative}) classifies each applied edit: descriptive prose
(``validates the token'') merely persists, while imperative prose (``should
validate the token'', ``Add retry logic''), an accepted unrealized plan node,
or retiring a feature that owns code each mint a \emph{directive} $d$ --- a
natural-language work order assembled from the feature's description and its
bound symbols, scoped to the files the feature owns. Directives are queued for
the user's \emph{live} coding-agent session rather than a headless model: the
session implements each directive, reflects the new code back through the
agent-side reflection API with $\delta = d$, and deletes the queue. Loop~A
then closes the cycle, with every resulting ledger entry carrying the causal
chain from the original document edit.

Verdict semantics are deliberately asymmetric for the destructive case
(line~\ref{line:retire}): accepting a machine-raised \textsf{retire} from the
inbox is \emph{detach-only} --- the feature is hidden and its chunks freed for
re-attribution, but no code-deletion directive is ever queued, because a retire
raised from transient drift may be a false positive and deleting code on a
single click is the most destructive failure mode. Only an explicit human
retire edit in the document text (or an agent proposal explicitly flagged to
delete code) queues removal work.

\begin{algorithm}[t]
\caption{Loop B --- tree $\to$ code realization (one pass)}
\label{alg:loop-b}
\begin{algorithmic}[1]
\Require tree text $T_{\mathrm{txt}}$; inbox verdicts $V$; intents $I$; store $(\T, \B, L)$
\Ensure updated $(\T, \B, L)$; directive queue $D$; refreshed $\sigma$
\State $\mathit{ops}_{\mathrm{user}} \gets \textsc{Diff}(T_{\mathrm{txt}}, \T)$ \Comment{snapshot BEFORE mutations (invariant 1)} \label{line:snapshot}
\State $D \gets \emptyset$
\For{$(e, \mathit{accept}) \in V$} \Comment{drain human verdicts on pending proposals}
    \If{$\mathit{accept}$}
        \State \textsc{Apply}$(e.\mathit{op}, \alpha{=}\textsf{human}, \mu{=}\pen)$; \textsc{Consume}$(e)$
        \If{$e.\mathit{op}$ is \textsf{retire} $\wedge\, \neg e.\mathit{op}.\mathit{deleteCode}$} \label{line:retire}
            \State \textsc{DetachAll}$(\B^{-1}(e.\mathit{op}.f))$ \Comment{detach-only: never queue code deletion}
        \ElsIf{\textsc{ImpliesCode}$(e.\mathit{op})$} $D \gets D \cup \{\textsc{Directive}(e.\mathit{op},\, \delta{=}e)\}$
        \EndIf
    \Else\ \textsc{Consume}$(e)$ \Comment{reject: delete the proposal}
    \EndIf
\EndFor
\For{$o \in \mathit{ops}_{\mathrm{user}}$} \Comment{direct edits are intentional: apply immediately}
    \State \textsc{Apply}$(o, \alpha, \mu \text{ from annotation channel (default \textsf{human}/\pen)})$
    \If{\textsc{ImpliesCode}$(o)$} $D \gets D \cup \{\textsc{Directive}(o, \delta{=}\ell_o)\}$ \label{line:imperative}
    \EndIf
\EndFor
\For{$i \in I$ with payload, fresh, unsatisfied} \Comment{doc-ahead suggestions: the loop applies them}
    \State \textsc{Apply}(\textsf{amend}$(i)$, $\alpha{=}i.\alpha$, $\mu{=}\suggest$, $\delta{=}i.\mathit{id}$)
    \If{\textsc{ImpliesCode}$(i)$} $D \gets D \cup \{\textsc{Directive}(i, \delta{=}i.\mathit{id})\}$
    \EndIf
\EndFor
\If{store mutated} \textsc{Render}$(\T) \to T_{\mathrm{txt}}$ \Comment{invariant 2: text catches up} \label{line:rerender}
\EndIf
\If{$D \neq \emptyset$}
    \State mint ids for $D$; queue $D$ for the live agent session; $H \gets H \cup \{d.f : d \in D\}$
    \State $\sigma \gets \awaiting$ \Comment{the agent implements, reflects with $\delta = d$, clears the queue}
\Else
    \State $\sigma \gets \textsc{RefreshState}()$
\EndIf
\State \Return $(\T, \B, L)$, $D$, $\sigma$
\end{algorithmic}
\end{algorithm}

\subsection{Conflict Resolution: the Ledger and Doc-Wins Holds}
\label{sec:conflicts}

Both loops route every decision through one classification table
(Table~\ref{tab:classify}): each detected change, from either side and by any
actor, maps to exactly one reaction. Three properties follow.

\paragraph{Provenance.} Every mutation is a ledger entry
(Definition~\ref{def:ledger}). The actor/mode pair separates \emph{who} changed
something from \emph{how committed} the change is, which drives both review
routing (only $\mu = \suggest$ items need verdicts) and the interface's
authorship rendering; $\delta$ threads causality across the realize cycle, so
a code change surfaced back by Loop~A is grouped under the document edit that
caused it.

\paragraph{Doc-wins holds.} Conflict resolution is a deterministic priority
rule rather than a merge heuristic. While $f \in H$ (a suggestion or queued
directive is pending on $f$), code-side \emph{intent} operations on $f$ ---
\textsf{amend}, \textsf{retire}, \textsf{move} --- are suppressed:
\begin{equation}
\label{eq:docwins}
\mathrm{suppress}(o, H) \;\iff\;
o.\mathit{kind} \in \{\textsf{amend}, \textsf{retire}, \textsf{move}\}
\;\wedge\; o.f \in H .
\end{equation}
Binding maintenance (\textsf{attach}/\textsf{detach}/\textsf{refresh}) is never
suppressed --- bindings are attribution, not intent, and must stay correct
regardless of who is mid-edit. Holds release when the suggestion is drained or
the directive queue completes, with a staleness backstop so an abandoned
suggestion cannot hold a feature indefinitely. The rule is the
tree-is-authoritative analogue of last-writer-wins: the side that expresses
\emph{intent} (the document) always outranks the side that merely
\emph{observes} (the code reflector).

\paragraph{Verdict grammar.} Who resolves what is structural:
accept/reject is exclusively a \emph{human} verb on \emph{code-ahead} items
(machine proposals), while a \emph{doc-ahead} suggestion is applied by the
machine side and offers its human author only \emph{withdraw}. No state is
double-resolvable, which eliminates accept/revert races by construction.

\begin{table}[t]
\centering
\caption{The change-classification table. Every change, from either side and by
any actor, maps to one reaction. $\auto$ reactions apply immediately;
$\suggest$ reactions await a human verdict; \textbf{D} marks reactions that
mint a realization directive.}
\label{tab:classify}
\small
\setlength{\tabcolsep}{4.5pt}
\begin{tabular}{@{}cllll@{}}
\toprule
\# & Side & Detected change & Origin & Reaction \\
\midrule
1 & code & modified bound chunk & any & \textsf{refresh} binding ($\auto$) \\
2 & code & removed bound chunk & any & \textsf{detach} ($\auto$); emptied $\wedge\, f \notin H \Rightarrow$ \textsf{retire} ($\suggest$, auth.\ pass) \\
3 & code & relocated chunk (Def.~\ref{def:reloc}) & any & deterministic \textsf{attach} to prior feature ($\auto$) \\
4 & code & added unbound chunk & any & LLM placement ($\auto$/$\suggest$); fallback $\mathrm{nbf}$ / \textsf{add} ($\suggest$) \\
5 & code & in-place modify, realized + prose & any & stale-prose \textsf{amend} (small $\to \auto$, large $\to \suggest$) \\
6 & code & any change during realize epoch & agent & rows 1--5 with $\delta = $ directive id \\
\midrule
7 & doc & descriptive amend / title edit & human, $\pen$ & apply immediately; no directive \\
8 & doc & imperative amend / plan add / retire-with-code & human, $\pen$ & apply + mint directive (\textbf{D}) \\
9 & doc & any edit & human, $\suggest$ & pending intent (hold); applied by Loop B drain $\to$ rows 7/8 \\
10 & doc & structural op via reflection API & agent & pending proposal ($\suggest$) \\
11 & doc & safe op via reflection API & agent & auto-apply; ledger $\alpha = \textsf{agent}$, $\mu = \auto$ \\
\midrule
12 & verdict & accept / reject & human & apply+consume / delete (\textsf{retire} accept is detach-only) \\
13 & code & drift on held feature ($f \in H$) & any & bindings maintained; \textsf{amend}/\textsf{retire}/\textsf{move} suppressed (Eq.~\ref{eq:docwins}) \\
\bottomrule
\end{tabular}
\end{table}

\subsection{Bootstrap}
\label{sec:bootstrap}

The initial tree is built by two scoped passes rather than one global
clustering. A \emph{per-file} pass issues one model call per source file,
seeing only that file's chunks and their incident graph edges --- making
cross-file ``junk drawer'' features structurally impossible --- and a
deterministic coverage step folds any chunk the model missed into the file's
largest feature. An \emph{organization} pass then groups the resulting
file-level features under a small number of thematic parents, using
feature-to-feature coupling aggregated from $\G$ as evidence. Both passes emit
ordinary ledger events through the same application seam as the loops, so
bootstrap output is curated (and corrected) with exactly the same verdict
machinery as any later change.
 from the ICLR template main file.
% Requires (all standard in the ICLR template ecosystem):
%   amsmath, amssymb, amsthm, booktabs, algorithm, algpseudocode
% plus the \newtheorem{definition}{Definition} declaration and the
% `problembox` environment + notation macros defined in main.tex.
% ------------------------------------------------------------------------------

\section{The \codoc{} System}
\label{sec:method}

\codoc{} maintains a human-authored \emph{feature tree} --- a hierarchy of named
units of intent --- bidirectionally synchronized with the code of a repository.
Unlike repository-level code graphs built \emph{from} code as retrieval context
for agents \citep{repograph2025,locagent2025,graphcoder2024,codexgraph2024}, the
tree is first-class authored intent: it is never regenerated from code, and the
attribution of code to features is a secondary, continuously repaired index.
Synchronization is asymmetric --- code changes are \emph{reflected} into the tree
(Loop~A, \S\ref{sec:loop-a}), while tree edits are \emph{realized} into code by a
coding agent (Loop~B, \S\ref{sec:loop-b}) --- and conflict resolution is
deterministic: pending document intent always wins over code-side observations
(\S\ref{sec:conflicts}).

\begin{problembox}{Problem statement}
Given a repository whose indexed chunk set is $\C$ and a human-authored feature
tree $\T = (\F, \prec)$, maintain (i)~a binding map
$\B \colon \C \to \F \cup \{\bot\}$ attributing every chunk to at most one
feature, and (ii)~the invariant that $\T$ expresses current human intent over
the code, under concurrent edits to either side by humans and autonomous
agents. Every mutation must be attributable (who, with what commitment, caused
by what), and divergence must resolve to a synchronization state
$\sigma = \synced$ without losing authored intent.
\end{problembox}

\subsection{Preliminaries}
\label{sec:prelim}

\begin{definition}[Chunk set]
\label{def:chunks}
The repository is indexed into a set of AST-level chunks
$\C = \{c_1, \dots, c_n\}$, where each chunk
$c = (\mathit{file},\, \mathit{sym},\, h^{\mathrm{tok}},\, h^{\mathrm{ast}})$
carries its file path, its qualified symbol path within the file, a
\emph{content hash} $h^{\mathrm{tok}}$ over the comment- and
whitespace-normalized token stream, and an \emph{AST-shape hash}
$h^{\mathrm{ast}}$ over the token-erased tree structure. The pair
$(\mathit{file}, \mathit{sym})$ is the chunk's identity key; the two hashes are
its relocation signals (Definition~\ref{def:reloc}).
\end{definition}

\begin{definition}[Feature tree]
\label{def:tree}
The feature tree is $\T = (\F, \prec)$, where $\F$ is a finite set of features
$f = (\mathit{id},\, \mathit{title},\, \mathit{desc},\, \mathit{realized})$ and
$\prec\, \subseteq \F \times \F$ is a forest-shaped parent relation. Titles and
descriptions are free prose authored by humans (or proposed by agents and
accepted by humans); $\mathit{realized} \in \{0,1\}$ marks whether the feature
is backed by code yet ($\mathit{realized}=0$ denotes an accepted plan
placeholder, flipped to $1$ by the first binding).
\end{definition}

\begin{definition}[Binding map]
\label{def:binding}
Attribution is a partial function $\B \colon \C \to \F \cup \{\bot\}$: each
chunk binds to at most one feature (enforced as a uniqueness constraint on the
chunk identity key), while one feature typically owns many chunks across many
files. Each binding stores the bound chunk's $h^{\mathrm{tok}}$ as its
\emph{fingerprint}, making staleness locally detectable. We write
$\B^{-1}(f) \subseteq \C$ for the chunks a feature owns and define coverage
$\mathrm{cov}(\B) = \lvert\{c : \B(c) \neq \bot\}\rvert / \lvert\C\rvert$.
\end{definition}

\begin{definition}[Code graph]
\label{def:graph}
Following the heterogeneous-graph convention of repository-level code
representations \citep{locagent2025,graphcoder2024}, the dependency structure
of the code is $\G = (\V, \E, \A, \R)$ with node-type map
$\tau \colon \V \to \A$, $\A = \{\textit{file}, \textit{class},
\textit{function}\}$, and edge-type map $\varphi \colon \E \to \R$,
$\R = \{\textit{contain}, \textit{call}, \textit{import}, \textit{inherit}\}$.
$\G$ is a derived cache, rebuildable from $\C$ at any time. The binding map
lifts $\G$ to a feature-coupling graph and defines the placement heuristic
\begin{equation}
\label{eq:neighbor}
\mathrm{nbf}(c) \;=\; \argmax_{f \in \F}\;
\bigl\lvert \{\, (v, v') \in \E \;:\; v' = \mathrm{sym}(c),\;
\B(\mathrm{chunk}(v)) = f \,\} \bigr\rvert ,
\end{equation}
the feature owning the most call/import edges incident to a new chunk $c$ ---
the deterministic fallback of Loop~A's coverage net
(Algorithm~\ref{alg:loop-a}, lines~\ref{line:covnet-start}--\ref{line:covnet-end}).
\end{definition}

\begin{definition}[Changeset]
\label{def:changeset}
A pass over the index yields $\Delta = (\Delta^{+}, \Delta^{-}, \Delta^{\sim})$:
chunks \emph{added} (present, unbound), \emph{removed} (a binding whose key has
left the index), and \emph{modified} (bound, but
$h^{\mathrm{tok}} \neq$ fingerprint). \codoc{} computes $\Delta$ in two modes:
a \emph{temporal} diff of consecutive index snapshots (cheap, runs on every
save) and a \emph{state-based} reconciliation that compares the index directly
against $\B$ (idempotent and self-healing; the authority pass). Only the latter
may retire features or amend prose, so destructive judgments are never made
from a transient mid-edit view.
\end{definition}

\begin{definition}[Relocation]
\label{def:reloc}
For $(r, a) \in \Delta^{-} \times \Delta^{+}$: $a$ is a \emph{move} of $r$ iff
$h^{\mathrm{tok}}_a = h^{\mathrm{tok}}_r$ (identical content at a new key), and
a \emph{rename} of $r$ iff $h^{\mathrm{ast}}_a = h^{\mathrm{ast}}_r$ uniquely
within the same file (same shape, new name; $1{:}1$ matches only). Relocations
re-attach attribution deterministically --- no model call, no risk of dropped
or duplicated attribution.
\end{definition}

\begin{definition}[Change ledger]
\label{def:ledger}
Every mutation of $(\T, \B)$ is an event in an append-only ledger $L$. A ledger
entry $\ell = (\mathit{op},\, f,\, \alpha,\, \mu,\, \delta)$ records the
operation, its target, the \emph{actor}
$\alpha \in \{\textsf{human}\} \cup \mathit{Agents} \cup \{\textsf{loop}\}$,
the \emph{mode} $\mu \in \{\pen, \suggest, \auto\}$ (committed edit, reviewable
proposal, or mechanical maintenance), and a \emph{causality link} $\delta$ ---
the directive, suggestion, or event id this change implements ($\bot$ if
spontaneous). A \emph{proposal} is an unapplied event; accepting it applies the
operation and consumes the event.
\end{definition}

\begin{definition}[Hold set and synchronization state]
\label{def:hold}
$H \subseteq \F$ is the set of features with pending document-ahead intent: a
live suggestion not yet drained, or a queued directive not yet implemented.
The synchronization state is
$\sigma \in \{\synced, \codedrift, \treedirty, \awaiting, \realizing\}$,
derived (not stored): pending code-side proposals put $\sigma = \codedrift$, a
non-empty directive queue forces $\sigma = \awaiting$ (a later code-side pass
cannot silently clobber queued work), and an active implementation pass reports
$\sigma = \realizing$.
\end{definition}

Table~\ref{tab:notation} summarizes the notation.

\begin{table}[t]
\centering
\caption{Notation.}
\label{tab:notation}
\small
\begin{tabular}{@{}ll@{}}
\toprule
Symbol & Meaning \\
\midrule
$\C$ & chunk set; $c = (\mathit{file}, \mathit{sym}, h^{\mathrm{tok}}, h^{\mathrm{ast}})$ \\
$\T = (\F, \prec)$ & feature tree (features + forest parent relation) \\
$\B \colon \C \to \F \cup \{\bot\}$ & binding map (each chunk owned by $\leq 1$ feature) \\
$\G = (\V, \E, \A, \R)$ & heterogeneous code dependency graph (derived cache) \\
$\Delta = (\Delta^{+}, \Delta^{-}, \Delta^{\sim})$ & changeset: added / removed / modified chunks \\
$h^{\mathrm{tok}}, h^{\mathrm{ast}}$ & content hash (move-invariant), AST-shape hash (rename-invariant) \\
$L$; $\ell = (\mathit{op}, f, \alpha, \mu, \delta)$ & append-only ledger; entry = op, target, actor, mode, cause \\
$H \subseteq \F$ & hold set (features with pending doc-ahead intent) \\
$\sigma$ & synchronization state \\
$\mathrm{nbf}(c)$ & graph-neighbor feature of a chunk (Eq.~\ref{eq:neighbor}) \\
$D$; $d$ & directive queue; a single realization directive \\
\bottomrule
\end{tabular}
\end{table}

\subsection{Loop A: Reflecting Code into the Tree}
\label{sec:loop-a}

Loop~A (Algorithm~\ref{alg:loop-a}) closes the gap from code to intent. Its
design principle is \emph{determinism first}: every change that can be resolved
mechanically --- refreshing fingerprints, detaching deleted chunks, carrying
attribution across moves and renames --- is resolved without a model
(lines~\ref{line:auto-start}--\ref{line:reloc-end}). A single LLM pass
(line~\ref{line:llm}) handles only the judgment calls that remain: placing
genuinely new code, repairing features that lost their last binding, and (in
the authority pass) refreshing prose whose code changed underneath it. The
pass sees the residual changeset, a file-locality subtree of $\T$, and
\emph{every} feature title; supplying the full title set to one call is what
prevents duplicate features without any post-hoc deduplication machinery.

Model output is not trusted with destructive authority. Proposed operations
are split by safety: binding maintenance and small description amendments
(bounded by an edit-distance ratio) auto-apply with $\mu = \auto$; structural
operations (add / move / retire) become unapplied proposal events with
$\mu = \suggest$ that a human accepts or rejects. Operations targeting held
features are dropped (line~\ref{line:hold}; \S\ref{sec:conflicts}), and a
\emph{coverage net}
(lines~\ref{line:covnet-start}--\ref{line:covnet-end}) guarantees no added
chunk is silently dropped: anything the model failed to place attaches to its
graph-neighbor feature $\mathrm{nbf}(c)$ or surfaces as an explicit
\textsf{add} proposal. Finally, stale-proposal garbage collection keeps the
pass idempotent: a pending \textsf{add} whose chunks have since bound
elsewhere, or a pending \textsf{retire} whose feature regained code, is
dropped, so repeated passes converge to $\sigma = \synced$ instead of
accumulating false positives.

\begin{algorithm}[t]
\caption{Loop A --- code $\to$ tree reflection (one pass)}
\label{alg:loop-a}
\begin{algorithmic}[1]
\Require index snapshot of $\C$; store $(\T, \B, L)$; hold set $H$; authority flag $\mathit{auth}$
\Ensure updated $(\T, \B, L)$; refreshed $\sigma$
\State $\Delta \gets \textsc{Changeset}(\C, \B)$ \Comment{temporal diff, or state-based iff $\mathit{auth}$}
\State \textsc{GcStaleProposals}$(L, \Delta)$ \Comment{drop superseded \textsf{add}/\textsf{retire} proposals}
\For{$c \in \Delta^{\sim}$} \textsc{Apply}(\textsf{refresh}$(c)$, $\mu{=}\auto$) \EndFor \label{line:auto-start}
\For{$c \in \Delta^{-}$} \textsc{Apply}(\textsf{detach}$(c)$, $\mu{=}\auto$) \EndFor
\For{$(r, a) \in \textsc{Relocations}(\Delta^{-}, \Delta^{+})$} \Comment{Def.~\ref{def:reloc}: $h^{\mathrm{tok}}$ move / $h^{\mathrm{ast}}$ rename}
    \State \textsc{Apply}(\textsf{attach}$(a, \B(r))$, $\mu{=}\auto$) \label{line:reloc-end}
\EndFor
\State $U \gets \{c \in \Delta^{+} : \B(c) = \bot\}$;\quad
       $E \gets \{f : \B^{-1}(f) = \emptyset \wedge f \notin H\}$ \Comment{unbound adds; emptied features}
\If{$U \neq \emptyset \,\vee\, E \neq \emptyset \,\vee\, (\mathit{auth} \wedge \textsc{StaleProse}(\Delta^{\sim}, H))$}
    \State $\mathit{ops} \gets \textsc{LLM}\bigl(\Delta,\ \textsc{Subtree}(\T, \Delta),\ \textsc{Titles}(\T),\ \G\bigr)$ \Comment{one call per pass} \label{line:llm}
    \State $\mathit{ops} \gets \{o \in \mathit{ops} : \neg\,\textsc{SuppressedByHold}(o, H)\}$ \label{line:hold}
    \For{$o \in \mathit{ops}$}
        \If{\textsc{Safe}$(o)$} \textsc{Apply}$(o, \mu{=}\auto)$
        \Else\ \textsc{Propose}$(o, \mu{=}\suggest)$ \Comment{unapplied event; human verdict pending}
        \EndIf
    \EndFor
\EndIf
\For{$c \in U$ with $\B(c) = \bot$ still} \label{line:covnet-start} \Comment{coverage net: nothing is dropped}
    \If{$\mathrm{nbf}(c) \neq \bot$} \textsc{Apply}(\textsf{attach}$(c, \mathrm{nbf}(c))$, $\mu{=}\auto$)
    \Else\ \textsc{Propose}(\textsf{add}$(c)$, $\mu{=}\suggest$)
    \EndIf
\EndFor \label{line:covnet-end}
\State \Return $(\T, \B, L)$, \textsc{RefreshState}$()$
\end{algorithmic}
\end{algorithm}

Two authority levels share this pipeline. The \emph{temporal} pass (run on
every save) computes $\Delta$ from consecutive snapshots and may never retire a
feature or amend prose --- a feature that looks emptied mid-edit usually
rebinds on the next save. The \emph{authority} pass (run at agent-session
close, on daemon startup, and from explicit sync) computes $\Delta$ state-based
against $\B$, is idempotent, recovers missed cycles, and is the only pass
allowed to raise \textsf{retire} and stale-prose \textsf{amend} proposals.

\subsection{Loop B: Realizing Tree Edits into Code}
\label{sec:loop-b}

Loop~B (Algorithm~\ref{alg:loop-b}) closes the opposite gap. Edits to the
rendered tree document are parsed and diffed against the store; verdicts on
pending proposals are drained from the review inbox; and live document-ahead
suggestions are applied by the loop itself (the human's only verb on their own
suggestion is \emph{withdraw} --- application is the agent side's
acknowledgment, stamped $\mu = \suggest$ with $\delta$ = the suggestion id).

Two ordering invariants keep text and store coherent
(lines~\ref{line:snapshot} and~\ref{line:rerender}). The text diff is
snapshotted \emph{before} any store mutation: verdict accepts move the store
ahead of the on-disk text, and diffing afterwards would misread the stale text
as a human edit that reverts the accepted change. Symmetrically, any pass that
mutated the store re-renders the text before returning, so the next pass diffs
a caught-up document.

Not every tree edit means ``write code.'' The \emph{imperative gate}
(line~\ref{line:imperative}) classifies each applied edit: descriptive prose
(``validates the token'') merely persists, while imperative prose (``should
validate the token'', ``Add retry logic''), an accepted unrealized plan node,
or retiring a feature that owns code each mint a \emph{directive} $d$ --- a
natural-language work order assembled from the feature's description and its
bound symbols, scoped to the files the feature owns. Directives are queued for
the user's \emph{live} coding-agent session rather than a headless model: the
session implements each directive, reflects the new code back through the
agent-side reflection API with $\delta = d$, and deletes the queue. Loop~A
then closes the cycle, with every resulting ledger entry carrying the causal
chain from the original document edit.

Verdict semantics are deliberately asymmetric for the destructive case
(line~\ref{line:retire}): accepting a machine-raised \textsf{retire} from the
inbox is \emph{detach-only} --- the feature is hidden and its chunks freed for
re-attribution, but no code-deletion directive is ever queued, because a retire
raised from transient drift may be a false positive and deleting code on a
single click is the most destructive failure mode. Only an explicit human
retire edit in the document text (or an agent proposal explicitly flagged to
delete code) queues removal work.

\begin{algorithm}[t]
\caption{Loop B --- tree $\to$ code realization (one pass)}
\label{alg:loop-b}
\begin{algorithmic}[1]
\Require tree text $T_{\mathrm{txt}}$; inbox verdicts $V$; intents $I$; store $(\T, \B, L)$
\Ensure updated $(\T, \B, L)$; directive queue $D$; refreshed $\sigma$
\State $\mathit{ops}_{\mathrm{user}} \gets \textsc{Diff}(T_{\mathrm{txt}}, \T)$ \Comment{snapshot BEFORE mutations (invariant 1)} \label{line:snapshot}
\State $D \gets \emptyset$
\For{$(e, \mathit{accept}) \in V$} \Comment{drain human verdicts on pending proposals}
    \If{$\mathit{accept}$}
        \State \textsc{Apply}$(e.\mathit{op}, \alpha{=}\textsf{human}, \mu{=}\pen)$; \textsc{Consume}$(e)$
        \If{$e.\mathit{op}$ is \textsf{retire} $\wedge\, \neg e.\mathit{op}.\mathit{deleteCode}$} \label{line:retire}
            \State \textsc{DetachAll}$(\B^{-1}(e.\mathit{op}.f))$ \Comment{detach-only: never queue code deletion}
        \ElsIf{\textsc{ImpliesCode}$(e.\mathit{op})$} $D \gets D \cup \{\textsc{Directive}(e.\mathit{op},\, \delta{=}e)\}$
        \EndIf
    \Else\ \textsc{Consume}$(e)$ \Comment{reject: delete the proposal}
    \EndIf
\EndFor
\For{$o \in \mathit{ops}_{\mathrm{user}}$} \Comment{direct edits are intentional: apply immediately}
    \State \textsc{Apply}$(o, \alpha, \mu \text{ from annotation channel (default \textsf{human}/\pen)})$
    \If{\textsc{ImpliesCode}$(o)$} $D \gets D \cup \{\textsc{Directive}(o, \delta{=}\ell_o)\}$ \label{line:imperative}
    \EndIf
\EndFor
\For{$i \in I$ with payload, fresh, unsatisfied} \Comment{doc-ahead suggestions: the loop applies them}
    \State \textsc{Apply}(\textsf{amend}$(i)$, $\alpha{=}i.\alpha$, $\mu{=}\suggest$, $\delta{=}i.\mathit{id}$)
    \If{\textsc{ImpliesCode}$(i)$} $D \gets D \cup \{\textsc{Directive}(i, \delta{=}i.\mathit{id})\}$
    \EndIf
\EndFor
\If{store mutated} \textsc{Render}$(\T) \to T_{\mathrm{txt}}$ \Comment{invariant 2: text catches up} \label{line:rerender}
\EndIf
\If{$D \neq \emptyset$}
    \State mint ids for $D$; queue $D$ for the live agent session; $H \gets H \cup \{d.f : d \in D\}$
    \State $\sigma \gets \awaiting$ \Comment{the agent implements, reflects with $\delta = d$, clears the queue}
\Else
    \State $\sigma \gets \textsc{RefreshState}()$
\EndIf
\State \Return $(\T, \B, L)$, $D$, $\sigma$
\end{algorithmic}
\end{algorithm}

\subsection{Conflict Resolution: the Ledger and Doc-Wins Holds}
\label{sec:conflicts}

Both loops route every decision through one classification table
(Table~\ref{tab:classify}): each detected change, from either side and by any
actor, maps to exactly one reaction. Three properties follow.

\paragraph{Provenance.} Every mutation is a ledger entry
(Definition~\ref{def:ledger}). The actor/mode pair separates \emph{who} changed
something from \emph{how committed} the change is, which drives both review
routing (only $\mu = \suggest$ items need verdicts) and the interface's
authorship rendering; $\delta$ threads causality across the realize cycle, so
a code change surfaced back by Loop~A is grouped under the document edit that
caused it.

\paragraph{Doc-wins holds.} Conflict resolution is a deterministic priority
rule rather than a merge heuristic. While $f \in H$ (a suggestion or queued
directive is pending on $f$), code-side \emph{intent} operations on $f$ ---
\textsf{amend}, \textsf{retire}, \textsf{move} --- are suppressed:
\begin{equation}
\label{eq:docwins}
\mathrm{suppress}(o, H) \;\iff\;
o.\mathit{kind} \in \{\textsf{amend}, \textsf{retire}, \textsf{move}\}
\;\wedge\; o.f \in H .
\end{equation}
Binding maintenance (\textsf{attach}/\textsf{detach}/\textsf{refresh}) is never
suppressed --- bindings are attribution, not intent, and must stay correct
regardless of who is mid-edit. Holds release when the suggestion is drained or
the directive queue completes, with a staleness backstop so an abandoned
suggestion cannot hold a feature indefinitely. The rule is the
tree-is-authoritative analogue of last-writer-wins: the side that expresses
\emph{intent} (the document) always outranks the side that merely
\emph{observes} (the code reflector).

\paragraph{Verdict grammar.} Who resolves what is structural:
accept/reject is exclusively a \emph{human} verb on \emph{code-ahead} items
(machine proposals), while a \emph{doc-ahead} suggestion is applied by the
machine side and offers its human author only \emph{withdraw}. No state is
double-resolvable, which eliminates accept/revert races by construction.

\begin{table}[t]
\centering
\caption{The change-classification table. Every change, from either side and by
any actor, maps to one reaction. $\auto$ reactions apply immediately;
$\suggest$ reactions await a human verdict; \textbf{D} marks reactions that
mint a realization directive.}
\label{tab:classify}
\small
\setlength{\tabcolsep}{4.5pt}
\begin{tabular}{@{}cllll@{}}
\toprule
\# & Side & Detected change & Origin & Reaction \\
\midrule
1 & code & modified bound chunk & any & \textsf{refresh} binding ($\auto$) \\
2 & code & removed bound chunk & any & \textsf{detach} ($\auto$); emptied $\wedge\, f \notin H \Rightarrow$ \textsf{retire} ($\suggest$, auth.\ pass) \\
3 & code & relocated chunk (Def.~\ref{def:reloc}) & any & deterministic \textsf{attach} to prior feature ($\auto$) \\
4 & code & added unbound chunk & any & LLM placement ($\auto$/$\suggest$); fallback $\mathrm{nbf}$ / \textsf{add} ($\suggest$) \\
5 & code & in-place modify, realized + prose & any & stale-prose \textsf{amend} (small $\to \auto$, large $\to \suggest$) \\
6 & code & any change during realize epoch & agent & rows 1--5 with $\delta = $ directive id \\
\midrule
7 & doc & descriptive amend / title edit & human, $\pen$ & apply immediately; no directive \\
8 & doc & imperative amend / plan add / retire-with-code & human, $\pen$ & apply + mint directive (\textbf{D}) \\
9 & doc & any edit & human, $\suggest$ & pending intent (hold); applied by Loop B drain $\to$ rows 7/8 \\
10 & doc & structural op via reflection API & agent & pending proposal ($\suggest$) \\
11 & doc & safe op via reflection API & agent & auto-apply; ledger $\alpha = \textsf{agent}$, $\mu = \auto$ \\
\midrule
12 & verdict & accept / reject & human & apply+consume / delete (\textsf{retire} accept is detach-only) \\
13 & code & drift on held feature ($f \in H$) & any & bindings maintained; \textsf{amend}/\textsf{retire}/\textsf{move} suppressed (Eq.~\ref{eq:docwins}) \\
\bottomrule
\end{tabular}
\end{table}

\subsection{Bootstrap}
\label{sec:bootstrap}

The initial tree is built by two scoped passes rather than one global
clustering. A \emph{per-file} pass issues one model call per source file,
seeing only that file's chunks and their incident graph edges --- making
cross-file ``junk drawer'' features structurally impossible --- and a
deterministic coverage step folds any chunk the model missed into the file's
largest feature. An \emph{organization} pass then groups the resulting
file-level features under a small number of thematic parents, using
feature-to-feature coupling aggregated from $\G$ as evidence. Both passes emit
ordinary ledger events through the same application seam as the loops, so
bootstrap output is curated (and corrected) with exactly the same verdict
machinery as any later change.
 from the ICLR template main file.
% Requires (all standard in the ICLR template ecosystem):
%   amsmath, amssymb, amsthm, booktabs, algorithm, algpseudocode
% plus the \newtheorem{definition}{Definition} declaration and the
% `problembox` environment + notation macros defined in main.tex.
% ------------------------------------------------------------------------------

\section{The \codoc{} System}
\label{sec:method}

\codoc{} maintains a human-authored \emph{feature tree} --- a hierarchy of named
units of intent --- bidirectionally synchronized with the code of a repository.
Unlike repository-level code graphs built \emph{from} code as retrieval context
for agents \citep{repograph2025,locagent2025,graphcoder2024,codexgraph2024}, the
tree is first-class authored intent: it is never regenerated from code, and the
attribution of code to features is a secondary, continuously repaired index.
Synchronization is asymmetric --- code changes are \emph{reflected} into the tree
(Loop~A, \S\ref{sec:loop-a}), while tree edits are \emph{realized} into code by a
coding agent (Loop~B, \S\ref{sec:loop-b}) --- and conflict resolution is
deterministic: pending document intent always wins over code-side observations
(\S\ref{sec:conflicts}).

\begin{problembox}{Problem statement}
Given a repository whose indexed chunk set is $\C$ and a human-authored feature
tree $\T = (\F, \prec)$, maintain (i)~a binding map
$\B \colon \C \to \F \cup \{\bot\}$ attributing every chunk to at most one
feature, and (ii)~the invariant that $\T$ expresses current human intent over
the code, under concurrent edits to either side by humans and autonomous
agents. Every mutation must be attributable (who, with what commitment, caused
by what), and divergence must resolve to a synchronization state
$\sigma = \synced$ without losing authored intent.
\end{problembox}

\subsection{Preliminaries}
\label{sec:prelim}

\begin{definition}[Chunk set]
\label{def:chunks}
The repository is indexed into a set of AST-level chunks
$\C = \{c_1, \dots, c_n\}$, where each chunk
$c = (\mathit{file},\, \mathit{sym},\, h^{\mathrm{tok}},\, h^{\mathrm{ast}})$
carries its file path, its qualified symbol path within the file, a
\emph{content hash} $h^{\mathrm{tok}}$ over the comment- and
whitespace-normalized token stream, and an \emph{AST-shape hash}
$h^{\mathrm{ast}}$ over the token-erased tree structure. The pair
$(\mathit{file}, \mathit{sym})$ is the chunk's identity key; the two hashes are
its relocation signals (Definition~\ref{def:reloc}).
\end{definition}

\begin{definition}[Feature tree]
\label{def:tree}
The feature tree is $\T = (\F, \prec)$, where $\F$ is a finite set of features
$f = (\mathit{id},\, \mathit{title},\, \mathit{desc},\, \mathit{realized})$ and
$\prec\, \subseteq \F \times \F$ is a forest-shaped parent relation. Titles and
descriptions are free prose authored by humans (or proposed by agents and
accepted by humans); $\mathit{realized} \in \{0,1\}$ marks whether the feature
is backed by code yet ($\mathit{realized}=0$ denotes an accepted plan
placeholder, flipped to $1$ by the first binding).
\end{definition}

\begin{definition}[Binding map]
\label{def:binding}
Attribution is a partial function $\B \colon \C \to \F \cup \{\bot\}$: each
chunk binds to at most one feature (enforced as a uniqueness constraint on the
chunk identity key), while one feature typically owns many chunks across many
files. Each binding stores the bound chunk's $h^{\mathrm{tok}}$ as its
\emph{fingerprint}, making staleness locally detectable. We write
$\B^{-1}(f) \subseteq \C$ for the chunks a feature owns and define coverage
$\mathrm{cov}(\B) = \lvert\{c : \B(c) \neq \bot\}\rvert / \lvert\C\rvert$.
\end{definition}

\begin{definition}[Code graph]
\label{def:graph}
Following the heterogeneous-graph convention of repository-level code
representations \citep{locagent2025,graphcoder2024}, the dependency structure
of the code is $\G = (\V, \E, \A, \R)$ with node-type map
$\tau \colon \V \to \A$, $\A = \{\textit{file}, \textit{class},
\textit{function}\}$, and edge-type map $\varphi \colon \E \to \R$,
$\R = \{\textit{contain}, \textit{call}, \textit{import}, \textit{inherit}\}$.
$\G$ is a derived cache, rebuildable from $\C$ at any time. The binding map
lifts $\G$ to a feature-coupling graph and defines the placement heuristic
\begin{equation}
\label{eq:neighbor}
\mathrm{nbf}(c) \;=\; \argmax_{f \in \F}\;
\bigl\lvert \{\, (v, v') \in \E \;:\; v' = \mathrm{sym}(c),\;
\B(\mathrm{chunk}(v)) = f \,\} \bigr\rvert ,
\end{equation}
the feature owning the most call/import edges incident to a new chunk $c$ ---
the deterministic fallback of Loop~A's coverage net
(Algorithm~\ref{alg:loop-a}, lines~\ref{line:covnet-start}--\ref{line:covnet-end}).
\end{definition}

\begin{definition}[Changeset]
\label{def:changeset}
A pass over the index yields $\Delta = (\Delta^{+}, \Delta^{-}, \Delta^{\sim})$:
chunks \emph{added} (present, unbound), \emph{removed} (a binding whose key has
left the index), and \emph{modified} (bound, but
$h^{\mathrm{tok}} \neq$ fingerprint). \codoc{} computes $\Delta$ in two modes:
a \emph{temporal} diff of consecutive index snapshots (cheap, runs on every
save) and a \emph{state-based} reconciliation that compares the index directly
against $\B$ (idempotent and self-healing; the authority pass). Only the latter
may retire features or amend prose, so destructive judgments are never made
from a transient mid-edit view.
\end{definition}

\begin{definition}[Relocation]
\label{def:reloc}
For $(r, a) \in \Delta^{-} \times \Delta^{+}$: $a$ is a \emph{move} of $r$ iff
$h^{\mathrm{tok}}_a = h^{\mathrm{tok}}_r$ (identical content at a new key), and
a \emph{rename} of $r$ iff $h^{\mathrm{ast}}_a = h^{\mathrm{ast}}_r$ uniquely
within the same file (same shape, new name; $1{:}1$ matches only). Relocations
re-attach attribution deterministically --- no model call, no risk of dropped
or duplicated attribution.
\end{definition}

\begin{definition}[Change ledger]
\label{def:ledger}
Every mutation of $(\T, \B)$ is an event in an append-only ledger $L$. A ledger
entry $\ell = (\mathit{op},\, f,\, \alpha,\, \mu,\, \delta)$ records the
operation, its target, the \emph{actor}
$\alpha \in \{\textsf{human}\} \cup \mathit{Agents} \cup \{\textsf{loop}\}$,
the \emph{mode} $\mu \in \{\pen, \suggest, \auto\}$ (committed edit, reviewable
proposal, or mechanical maintenance), and a \emph{causality link} $\delta$ ---
the directive, suggestion, or event id this change implements ($\bot$ if
spontaneous). A \emph{proposal} is an unapplied event; accepting it applies the
operation and consumes the event.
\end{definition}

\begin{definition}[Hold set and synchronization state]
\label{def:hold}
$H \subseteq \F$ is the set of features with pending document-ahead intent: a
live suggestion not yet drained, or a queued directive not yet implemented.
The synchronization state is
$\sigma \in \{\synced, \codedrift, \treedirty, \awaiting, \realizing\}$,
derived (not stored): pending code-side proposals put $\sigma = \codedrift$, a
non-empty directive queue forces $\sigma = \awaiting$ (a later code-side pass
cannot silently clobber queued work), and an active implementation pass reports
$\sigma = \realizing$.
\end{definition}

Table~\ref{tab:notation} summarizes the notation.

\begin{table}[t]
\centering
\caption{Notation.}
\label{tab:notation}
\small
\begin{tabular}{@{}ll@{}}
\toprule
Symbol & Meaning \\
\midrule
$\C$ & chunk set; $c = (\mathit{file}, \mathit{sym}, h^{\mathrm{tok}}, h^{\mathrm{ast}})$ \\
$\T = (\F, \prec)$ & feature tree (features + forest parent relation) \\
$\B \colon \C \to \F \cup \{\bot\}$ & binding map (each chunk owned by $\leq 1$ feature) \\
$\G = (\V, \E, \A, \R)$ & heterogeneous code dependency graph (derived cache) \\
$\Delta = (\Delta^{+}, \Delta^{-}, \Delta^{\sim})$ & changeset: added / removed / modified chunks \\
$h^{\mathrm{tok}}, h^{\mathrm{ast}}$ & content hash (move-invariant), AST-shape hash (rename-invariant) \\
$L$; $\ell = (\mathit{op}, f, \alpha, \mu, \delta)$ & append-only ledger; entry = op, target, actor, mode, cause \\
$H \subseteq \F$ & hold set (features with pending doc-ahead intent) \\
$\sigma$ & synchronization state \\
$\mathrm{nbf}(c)$ & graph-neighbor feature of a chunk (Eq.~\ref{eq:neighbor}) \\
$D$; $d$ & directive queue; a single realization directive \\
\bottomrule
\end{tabular}
\end{table}

\subsection{Loop A: Reflecting Code into the Tree}
\label{sec:loop-a}

Loop~A (Algorithm~\ref{alg:loop-a}) closes the gap from code to intent. Its
design principle is \emph{determinism first}: every change that can be resolved
mechanically --- refreshing fingerprints, detaching deleted chunks, carrying
attribution across moves and renames --- is resolved without a model
(lines~\ref{line:auto-start}--\ref{line:reloc-end}). A single LLM pass
(line~\ref{line:llm}) handles only the judgment calls that remain: placing
genuinely new code, repairing features that lost their last binding, and (in
the authority pass) refreshing prose whose code changed underneath it. The
pass sees the residual changeset, a file-locality subtree of $\T$, and
\emph{every} feature title; supplying the full title set to one call is what
prevents duplicate features without any post-hoc deduplication machinery.

Model output is not trusted with destructive authority. Proposed operations
are split by safety: binding maintenance and small description amendments
(bounded by an edit-distance ratio) auto-apply with $\mu = \auto$; structural
operations (add / move / retire) become unapplied proposal events with
$\mu = \suggest$ that a human accepts or rejects. Operations targeting held
features are dropped (line~\ref{line:hold}; \S\ref{sec:conflicts}), and a
\emph{coverage net}
(lines~\ref{line:covnet-start}--\ref{line:covnet-end}) guarantees no added
chunk is silently dropped: anything the model failed to place attaches to its
graph-neighbor feature $\mathrm{nbf}(c)$ or surfaces as an explicit
\textsf{add} proposal. Finally, stale-proposal garbage collection keeps the
pass idempotent: a pending \textsf{add} whose chunks have since bound
elsewhere, or a pending \textsf{retire} whose feature regained code, is
dropped, so repeated passes converge to $\sigma = \synced$ instead of
accumulating false positives.

\begin{algorithm}[t]
\caption{Loop A --- code $\to$ tree reflection (one pass)}
\label{alg:loop-a}
\begin{algorithmic}[1]
\Require index snapshot of $\C$; store $(\T, \B, L)$; hold set $H$; authority flag $\mathit{auth}$
\Ensure updated $(\T, \B, L)$; refreshed $\sigma$
\State $\Delta \gets \textsc{Changeset}(\C, \B)$ \Comment{temporal diff, or state-based iff $\mathit{auth}$}
\State \textsc{GcStaleProposals}$(L, \Delta)$ \Comment{drop superseded \textsf{add}/\textsf{retire} proposals}
\For{$c \in \Delta^{\sim}$} \textsc{Apply}(\textsf{refresh}$(c)$, $\mu{=}\auto$) \EndFor \label{line:auto-start}
\For{$c \in \Delta^{-}$} \textsc{Apply}(\textsf{detach}$(c)$, $\mu{=}\auto$) \EndFor
\For{$(r, a) \in \textsc{Relocations}(\Delta^{-}, \Delta^{+})$} \Comment{Def.~\ref{def:reloc}: $h^{\mathrm{tok}}$ move / $h^{\mathrm{ast}}$ rename}
    \State \textsc{Apply}(\textsf{attach}$(a, \B(r))$, $\mu{=}\auto$) \label{line:reloc-end}
\EndFor
\State $U \gets \{c \in \Delta^{+} : \B(c) = \bot\}$;\quad
       $E \gets \{f : \B^{-1}(f) = \emptyset \wedge f \notin H\}$ \Comment{unbound adds; emptied features}
\If{$U \neq \emptyset \,\vee\, E \neq \emptyset \,\vee\, (\mathit{auth} \wedge \textsc{StaleProse}(\Delta^{\sim}, H))$}
    \State $\mathit{ops} \gets \textsc{LLM}\bigl(\Delta,\ \textsc{Subtree}(\T, \Delta),\ \textsc{Titles}(\T),\ \G\bigr)$ \Comment{one call per pass} \label{line:llm}
    \State $\mathit{ops} \gets \{o \in \mathit{ops} : \neg\,\textsc{SuppressedByHold}(o, H)\}$ \label{line:hold}
    \For{$o \in \mathit{ops}$}
        \If{\textsc{Safe}$(o)$} \textsc{Apply}$(o, \mu{=}\auto)$
        \Else\ \textsc{Propose}$(o, \mu{=}\suggest)$ \Comment{unapplied event; human verdict pending}
        \EndIf
    \EndFor
\EndIf
\For{$c \in U$ with $\B(c) = \bot$ still} \label{line:covnet-start} \Comment{coverage net: nothing is dropped}
    \If{$\mathrm{nbf}(c) \neq \bot$} \textsc{Apply}(\textsf{attach}$(c, \mathrm{nbf}(c))$, $\mu{=}\auto$)
    \Else\ \textsc{Propose}(\textsf{add}$(c)$, $\mu{=}\suggest$)
    \EndIf
\EndFor \label{line:covnet-end}
\State \Return $(\T, \B, L)$, \textsc{RefreshState}$()$
\end{algorithmic}
\end{algorithm}

Two authority levels share this pipeline. The \emph{temporal} pass (run on
every save) computes $\Delta$ from consecutive snapshots and may never retire a
feature or amend prose --- a feature that looks emptied mid-edit usually
rebinds on the next save. The \emph{authority} pass (run at agent-session
close, on daemon startup, and from explicit sync) computes $\Delta$ state-based
against $\B$, is idempotent, recovers missed cycles, and is the only pass
allowed to raise \textsf{retire} and stale-prose \textsf{amend} proposals.

\subsection{Loop B: Realizing Tree Edits into Code}
\label{sec:loop-b}

Loop~B (Algorithm~\ref{alg:loop-b}) closes the opposite gap. Edits to the
rendered tree document are parsed and diffed against the store; verdicts on
pending proposals are drained from the review inbox; and live document-ahead
suggestions are applied by the loop itself (the human's only verb on their own
suggestion is \emph{withdraw} --- application is the agent side's
acknowledgment, stamped $\mu = \suggest$ with $\delta$ = the suggestion id).

Two ordering invariants keep text and store coherent
(lines~\ref{line:snapshot} and~\ref{line:rerender}). The text diff is
snapshotted \emph{before} any store mutation: verdict accepts move the store
ahead of the on-disk text, and diffing afterwards would misread the stale text
as a human edit that reverts the accepted change. Symmetrically, any pass that
mutated the store re-renders the text before returning, so the next pass diffs
a caught-up document.

Not every tree edit means ``write code.'' The \emph{imperative gate}
(line~\ref{line:imperative}) classifies each applied edit: descriptive prose
(``validates the token'') merely persists, while imperative prose (``should
validate the token'', ``Add retry logic''), an accepted unrealized plan node,
or retiring a feature that owns code each mint a \emph{directive} $d$ --- a
natural-language work order assembled from the feature's description and its
bound symbols, scoped to the files the feature owns. Directives are queued for
the user's \emph{live} coding-agent session rather than a headless model: the
session implements each directive, reflects the new code back through the
agent-side reflection API with $\delta = d$, and deletes the queue. Loop~A
then closes the cycle, with every resulting ledger entry carrying the causal
chain from the original document edit.

Verdict semantics are deliberately asymmetric for the destructive case
(line~\ref{line:retire}): accepting a machine-raised \textsf{retire} from the
inbox is \emph{detach-only} --- the feature is hidden and its chunks freed for
re-attribution, but no code-deletion directive is ever queued, because a retire
raised from transient drift may be a false positive and deleting code on a
single click is the most destructive failure mode. Only an explicit human
retire edit in the document text (or an agent proposal explicitly flagged to
delete code) queues removal work.

\begin{algorithm}[t]
\caption{Loop B --- tree $\to$ code realization (one pass)}
\label{alg:loop-b}
\begin{algorithmic}[1]
\Require tree text $T_{\mathrm{txt}}$; inbox verdicts $V$; intents $I$; store $(\T, \B, L)$
\Ensure updated $(\T, \B, L)$; directive queue $D$; refreshed $\sigma$
\State $\mathit{ops}_{\mathrm{user}} \gets \textsc{Diff}(T_{\mathrm{txt}}, \T)$ \Comment{snapshot BEFORE mutations (invariant 1)} \label{line:snapshot}
\State $D \gets \emptyset$
\For{$(e, \mathit{accept}) \in V$} \Comment{drain human verdicts on pending proposals}
    \If{$\mathit{accept}$}
        \State \textsc{Apply}$(e.\mathit{op}, \alpha{=}\textsf{human}, \mu{=}\pen)$; \textsc{Consume}$(e)$
        \If{$e.\mathit{op}$ is \textsf{retire} $\wedge\, \neg e.\mathit{op}.\mathit{deleteCode}$} \label{line:retire}
            \State \textsc{DetachAll}$(\B^{-1}(e.\mathit{op}.f))$ \Comment{detach-only: never queue code deletion}
        \ElsIf{\textsc{ImpliesCode}$(e.\mathit{op})$} $D \gets D \cup \{\textsc{Directive}(e.\mathit{op},\, \delta{=}e)\}$
        \EndIf
    \Else\ \textsc{Consume}$(e)$ \Comment{reject: delete the proposal}
    \EndIf
\EndFor
\For{$o \in \mathit{ops}_{\mathrm{user}}$} \Comment{direct edits are intentional: apply immediately}
    \State \textsc{Apply}$(o, \alpha, \mu \text{ from annotation channel (default \textsf{human}/\pen)})$
    \If{\textsc{ImpliesCode}$(o)$} $D \gets D \cup \{\textsc{Directive}(o, \delta{=}\ell_o)\}$ \label{line:imperative}
    \EndIf
\EndFor
\For{$i \in I$ with payload, fresh, unsatisfied} \Comment{doc-ahead suggestions: the loop applies them}
    \State \textsc{Apply}(\textsf{amend}$(i)$, $\alpha{=}i.\alpha$, $\mu{=}\suggest$, $\delta{=}i.\mathit{id}$)
    \If{\textsc{ImpliesCode}$(i)$} $D \gets D \cup \{\textsc{Directive}(i, \delta{=}i.\mathit{id})\}$
    \EndIf
\EndFor
\If{store mutated} \textsc{Render}$(\T) \to T_{\mathrm{txt}}$ \Comment{invariant 2: text catches up} \label{line:rerender}
\EndIf
\If{$D \neq \emptyset$}
    \State mint ids for $D$; queue $D$ for the live agent session; $H \gets H \cup \{d.f : d \in D\}$
    \State $\sigma \gets \awaiting$ \Comment{the agent implements, reflects with $\delta = d$, clears the queue}
\Else
    \State $\sigma \gets \textsc{RefreshState}()$
\EndIf
\State \Return $(\T, \B, L)$, $D$, $\sigma$
\end{algorithmic}
\end{algorithm}

\subsection{Conflict Resolution: the Ledger and Doc-Wins Holds}
\label{sec:conflicts}

Both loops route every decision through one classification table
(Table~\ref{tab:classify}): each detected change, from either side and by any
actor, maps to exactly one reaction. Three properties follow.

\paragraph{Provenance.} Every mutation is a ledger entry
(Definition~\ref{def:ledger}). The actor/mode pair separates \emph{who} changed
something from \emph{how committed} the change is, which drives both review
routing (only $\mu = \suggest$ items need verdicts) and the interface's
authorship rendering; $\delta$ threads causality across the realize cycle, so
a code change surfaced back by Loop~A is grouped under the document edit that
caused it.

\paragraph{Doc-wins holds.} Conflict resolution is a deterministic priority
rule rather than a merge heuristic. While $f \in H$ (a suggestion or queued
directive is pending on $f$), code-side \emph{intent} operations on $f$ ---
\textsf{amend}, \textsf{retire}, \textsf{move} --- are suppressed:
\begin{equation}
\label{eq:docwins}
\mathrm{suppress}(o, H) \;\iff\;
o.\mathit{kind} \in \{\textsf{amend}, \textsf{retire}, \textsf{move}\}
\;\wedge\; o.f \in H .
\end{equation}
Binding maintenance (\textsf{attach}/\textsf{detach}/\textsf{refresh}) is never
suppressed --- bindings are attribution, not intent, and must stay correct
regardless of who is mid-edit. Holds release when the suggestion is drained or
the directive queue completes, with a staleness backstop so an abandoned
suggestion cannot hold a feature indefinitely. The rule is the
tree-is-authoritative analogue of last-writer-wins: the side that expresses
\emph{intent} (the document) always outranks the side that merely
\emph{observes} (the code reflector).

\paragraph{Verdict grammar.} Who resolves what is structural:
accept/reject is exclusively a \emph{human} verb on \emph{code-ahead} items
(machine proposals), while a \emph{doc-ahead} suggestion is applied by the
machine side and offers its human author only \emph{withdraw}. No state is
double-resolvable, which eliminates accept/revert races by construction.

\begin{table}[t]
\centering
\caption{The change-classification table. Every change, from either side and by
any actor, maps to one reaction. $\auto$ reactions apply immediately;
$\suggest$ reactions await a human verdict; \textbf{D} marks reactions that
mint a realization directive.}
\label{tab:classify}
\small
\setlength{\tabcolsep}{4.5pt}
\begin{tabular}{@{}cllll@{}}
\toprule
\# & Side & Detected change & Origin & Reaction \\
\midrule
1 & code & modified bound chunk & any & \textsf{refresh} binding ($\auto$) \\
2 & code & removed bound chunk & any & \textsf{detach} ($\auto$); emptied $\wedge\, f \notin H \Rightarrow$ \textsf{retire} ($\suggest$, auth.\ pass) \\
3 & code & relocated chunk (Def.~\ref{def:reloc}) & any & deterministic \textsf{attach} to prior feature ($\auto$) \\
4 & code & added unbound chunk & any & LLM placement ($\auto$/$\suggest$); fallback $\mathrm{nbf}$ / \textsf{add} ($\suggest$) \\
5 & code & in-place modify, realized + prose & any & stale-prose \textsf{amend} (small $\to \auto$, large $\to \suggest$) \\
6 & code & any change during realize epoch & agent & rows 1--5 with $\delta = $ directive id \\
\midrule
7 & doc & descriptive amend / title edit & human, $\pen$ & apply immediately; no directive \\
8 & doc & imperative amend / plan add / retire-with-code & human, $\pen$ & apply + mint directive (\textbf{D}) \\
9 & doc & any edit & human, $\suggest$ & pending intent (hold); applied by Loop B drain $\to$ rows 7/8 \\
10 & doc & structural op via reflection API & agent & pending proposal ($\suggest$) \\
11 & doc & safe op via reflection API & agent & auto-apply; ledger $\alpha = \textsf{agent}$, $\mu = \auto$ \\
\midrule
12 & verdict & accept / reject & human & apply+consume / delete (\textsf{retire} accept is detach-only) \\
13 & code & drift on held feature ($f \in H$) & any & bindings maintained; \textsf{amend}/\textsf{retire}/\textsf{move} suppressed (Eq.~\ref{eq:docwins}) \\
\bottomrule
\end{tabular}
\end{table}

\subsection{Bootstrap}
\label{sec:bootstrap}

The initial tree is built by two scoped passes rather than one global
clustering. A \emph{per-file} pass issues one model call per source file,
seeing only that file's chunks and their incident graph edges --- making
cross-file ``junk drawer'' features structurally impossible --- and a
deterministic coverage step folds any chunk the model missed into the file's
largest feature. An \emph{organization} pass then groups the resulting
file-level features under a small number of thematic parents, using
feature-to-feature coupling aggregated from $\G$ as evidence. Both passes emit
ordinary ledger events through the same application seam as the loops, so
bootstrap output is curated (and corrected) with exactly the same verdict
machinery as any later change.


\bibliographystyle{plainnat}
\bibliography{references}

\end{document}
